\documentclass[12pt,a4paper]{article}
\usepackage[utf8]{inputenc}
\usepackage[T1]{fontenc}
\usepackage[ngerman,english]{babel}
\usepackage{amsmath,amssymb,amsthm}
\usepackage{geometry}
\geometry{margin=2.5cm}
\usepackage{fancyhdr}
\usepackage{titlesec}

% Comprehensive macro block
\newcommand{\cO}{\mathcal{O}}
\newcommand{\cF}{\mathcal{F}}
\newcommand{\cE}{\mathcal{E}}
\newcommand{\cL}{\mathcal{L}}
\newcommand{\cH}{\mathcal{H}}
\newcommand{\cA}{\mathcal{A}}
\newcommand{\cB}{\mathcal{B}}
\newcommand{\cC}{\mathcal{C}}
\newcommand{\cD}{\mathcal{D}}
\newcommand{\cG}{\mathcal{G}}
\newcommand{\cI}{\mathcal{I}}
\newcommand{\cJ}{\mathcal{J}}
\newcommand{\cK}{\mathcal{K}}
\newcommand{\cM}{\mathcal{M}}
\newcommand{\cN}{\mathcal{N}}
\newcommand{\cP}{\mathcal{P}}
\newcommand{\cS}{\mathcal{S}}
\newcommand{\cT}{\mathcal{T}}
\newcommand{\cU}{\mathcal{U}}
\newcommand{\cV}{\mathcal{V}}
\newcommand{\cW}{\mathcal{W}}
\newcommand{\bA}{\mathbb{A}}
\newcommand{\bC}{\mathbb{C}}
\newcommand{\bF}{\mathbb{F}}
\newcommand{\bG}{\mathbb{G}}
\newcommand{\bN}{\mathbb{N}}
\newcommand{\bP}{\mathbb{P}}
\newcommand{\bQ}{\mathbb{Q}}
\newcommand{\bR}{\mathbb{R}}
\newcommand{\bZ}{\mathbb{Z}}
\newcommand{\fm}{\mathfrak{m}}
\newcommand{\fp}{\mathfrak{p}}
\newcommand{\fq}{\mathfrak{q}}
\newcommand{\fP}{\mathfrak{P}}
\newcommand{\fa}{\mathfrak{a}}
\DeclareMathOperator{\Gal}{Gal}
\DeclareMathOperator{\Aut}{Aut}
\DeclareMathOperator{\End}{End}
\DeclareMathOperator{\Hom}{Hom}
\DeclareMathOperator{\Ext}{Ext}
\DeclareMathOperator{\Tor}{Tor}
\DeclareMathOperator{\Pic}{Pic}
\DeclareMathOperator{\Ker}{Ker}
\let\Im\relax
\DeclareMathOperator{\Im}{Im}
\DeclareMathOperator{\codim}{codim}
\DeclareMathOperator{\supp}{supp}
\DeclareMathOperator{\Spec}{Spec}
\DeclareMathOperator{\Spf}{Spf}
\DeclareMathOperator{\Proj}{Proj}
\DeclareMathOperator{\Div}{Div}
\DeclareMathOperator{\ord}{ord}
\DeclareMathOperator{\card}{card}
\DeclareMathOperator{\rg}{rg}
\DeclareMathOperator{\tr}{tr}
\DeclareMathOperator{\Tr}{Tr}
\DeclareMathOperator{\Res}{Res}
\DeclareMathOperator{\Ind}{Ind}
\DeclareMathOperator{\Cor}{Cor}
\DeclareMathOperator{\Br}{Br}
\DeclareMathOperator{\Alb}{Alb}
\DeclareMathOperator{\Jac}{Jac}
\DeclareMathOperator{\Prym}{Prym}
\DeclareMathOperator{\Sym}{Sym}
\DeclareMathOperator{\Alt}{Alt}
\DeclareMathOperator{\Num}{Num}
\DeclareMathOperator{\GL}{GL}

% Custom operators for invariant theory
\DeclareMathOperator{\Defekt}{Defekt}

\pagestyle{fancy}
\fancyhf{}
\fancyhead[LE,RO]{\thepage}
\renewcommand{\headrulewidth}{0pt}

\begin{document}
\renewenvironment{noetherpaper}[4]{\newpage\section*{#1}\addcontentsline{toc}{section}{#1}\par\smallskip}{\par\medskip}


% ============================================
% PAPER 3: Pages 101--104
% ============================================

\selectlanguage{ngerman}

\begin{center}
\textbf{3. Zur Invariantentheorie der Formen von $n$ Variablen}\footnote{Vortrag, gehalten auf der Salzburger Naturforscherversammlung 1909.}

\smallskip
\textit{J. Ber. d. DMV 19 (1910), S. 101--104}
\end{center}

\smallskip

In der projektiven Invariantentheorie der Formen von $n$ Variabeln ist das Hauptproblem gelöst, die Endlichkeit des Formensystems bewiesen. Eine Reihe weiterer Fragen aber sind nicht behandelt, nämlich solche, die sich auf Methoden zur Erforschung des Formenzusammenhangs beziehen. Die Resultate, die ich in bezug auf solche Fragen erhalten habe, möchte ich kurz mitteilen, der Deutlichkeit halber aber die Fragestellungen erst im ternären Gebiet, wo sie bekannt sind, skizzieren.

Ich denke zuerst an die beiden Fundamentalsätze der symbolischen Methode, den Satz, daß sich alle invarianten Bildungen symbolisch darstellen lassen, und den zweiten, daß man alle zwischen Invarianten bestehenden Relationen erhält durch sukzessive Anwendung einer kleinen Anzahl von symbolischen Identitäten; daß also die symbolische Methode alle Relationen gibt, ohne aus dem Gebiet des Invarianten herauszutreten.\footnote{Study: \textit{Methoden zur Theorie der ternären Formen.} II. \S\ 6. (Leipzig, Teubner 1889.)}

Darauf aufbauend kommen dann die Prozesse zur Erzeugung der Formen, der symbolische Faltungsprozeß und die unsymbolischen Differentiationsprozesse, die Theorie der Reihenentwicklungen u.\,a. Hierzu gehörig möchte ich noch einen in meiner Dissertation\footnote{\textit{Über die Bildung des Formensystems der ternären biquadratischen Form.} \S\ 1. (Crelle's Journal Bd.~134. 1908.)} bewiesenen Satz nennen, daß alle Faltungen erzeugt werden können durch sukzessive Anwendung der zwei möglichen ``kogredienten'' Faltungen; darunter verstehe ich bei einem symbolischen Produkt $s,t,u,v$ die beiden Faltungen $(stu)$ und $(vx)$. Auf diesen Satz läßt sich die Theorie der Reduzenten und der Reduktion der Formensysteme aufbauen, analog wie im binären Gebiet, wie ich in meiner Dissertation gezeigt habe.

Wenn wir nun zum $n$-ären Gebiet übergehen, so gilt schon der erste Satz der symbolischen Methode nur in bedingtem Maße. Bekanntlich treten schon bei den quaternären Formen außer den Punkt- und Ebenenkoordinaten $x$ und $u$ Linienkoordinaten $p_{\varrho}$ auf, die Unterdeterminanten aus zwei Reihen von Punktkoordinaten. Analog haben wir im $n$-ären Gebiet Variablenreihen $p_{\varrho}$, deren einzelne Elemente die Unterdeterminanten aus $\varrho$ Reihen $x$ sind, und Symbolreihen $S^{\varrho}$, deren Elemente sich verhalten wie Unterdeterminanten aus $\varrho$ Reihen $s$, die den $u$ kogredient sind. Die symbolische Darstellung durch Determinanten gilt nun bekanntlich nur, wenn die Symbolreihen $S^{\varrho}$ in die Reihen $s$ aufgelöst und diese in verschiedene Determinanten verteilt werden; eine reale Bedeutung haben dann nur solche Aggregate von Determinanten, die von den $S^{\varrho}$ abhängen. Welche Determinantenaggregate das leisten, läßt sich aber nicht übersehen; und es gehen dadurch alle übrigen der im ternären Gebiete genannten Sätze verloren.

Ich möchte nun ein einfaches Prinzip angeben, um eine in den Reihen $S^{\varrho}$ explizite symbolische Darstellung zu gewinnen und damit auch die übrigen Sätze wieder zu erhalten. Es ist dies eine Anwendung des bekannten Satzes über die korrespondierenden Matrizen: ``Einer jeden Variablenreihe $p_{\varrho}$ läßt sich ein-eindeutig eine andere $q_{n-\varrho}$ zuordnen, deren einzelne Elemente Unterdeterminanten aus $n-\varrho$ Reihen $u$ sind, die mit den $\varrho$ Reihen $x$ durch Gleichungen $(u|x)=0$ verbunden sind.''

Entsprechend ist jeder Symbolreihe $S^{\varrho}$ eine andere $S^{n-\varrho}$ zugeordnet, die sich verhält wie aus $n-\varrho$ den $x$ kogredienten Reihen zusammengesetzt.

Der Satz läßt noch eine zweite, für Anwendungen geeignete Fassung zu: ``Einem jeden Matrizenprodukt $(v^{(1)} \cdots v^{(\varrho)} | y^{(1)} \cdots y^{(\varrho)})$,\footnote{D.h.\ der Summe über die Produkte entsprechender Determinanten, gebildet aus der Matrix der $\varrho$ Reihen $v$ und derjenigen der $\varrho$ Reihen $y$.} wo die Reihen der zweiten Matrix denjenigen der ersten kontragredient sind, ist eine Determinante zugeordnet; und umgekehrt läßt sich jede Determinante in ein Matrizenprodukt von vorgegebener Reihenanzahl $\varrho$ verwandeln.'' Damit erscheinen aber Determinante und Matrizenprodukt als vollständig äquivalent, und der erste Satz der symbolischen Methode nimmt die Form an: ``Alle invarianten Bildungen lassen sich durch Matrizenprodukte symbolisch darstellen.''

Aus zwei Symbolreihen $S^{\varrho}$, $T^{\tau}$ läßt sich nun mit Hilfe von Variablenreihen sehr leicht ein diese Symbolreihen explizit enthaltendes Matrizenprodukt bilden, das also eine invariante Bildung der Form $(S^{\varrho} | p_{\varrho})(T^{\tau} | p_{\tau})$ darstellt, nämlich das folgende:
\begin{equation}\tag{1}
(G_{n-\varrho-\lambda} S^{\varrho} T^{\tau} | p_{\varrho+\lambda} p_{\tau-\lambda}) \qquad (\lambda = 1, 2, \ldots \tau \text{ oder } n-\varrho).
\end{equation}

Wichtiger ist die Umkehrung: daß sich alle invarianten Bildungen explizit durch Matrizenprodukte darstellen lassen, daß also der obige Prozeß die Erweiterung des binären und ternären Faltungsprozesses bedeutet. Zum Beweise dieser Umkehrung ist es nötig, den zweiten Satz der symbolischen Methode auf das $n$-äre Gebiet zu übertragen, nämlich die Gesamtheit der unabhängigen unzerlegbaren Identitäten aufzustellen, bei denen alle Reihen $S^{\varrho}$, $p_{\varrho}$ als unzerlegbar betrachtet werden. Diese Identitäten ergeben sich aus den bekannten, nur Reihen $x$ und $u$ enthaltenden,\footnote{E.\ Pascal in \textit{Memorie d.\ R.\ Acc.\ d.\ Lincei.} (4) 5 (1888), p.~375.} durch Ersetzen des Produktsatzes für Determinanten durch ein System von Produktsätzen für Matrizen, und Zusammenziehen verschiedener Matrizenprodukte zu einem einzigen durch eben diese Produktsätze. Man gelangt so zu zwei dualistischen Formeln, von denen nur eine angeführt sei:
\begin{equation}\tag{2}
\sum_{i=1}^{k} (S^{\varrho_{1}} \cdots S^{\varrho_{k}} | p_{\varrho} x^{(i)}) (S^{\varrho} | x^{(i)}) = (S^{\varrho_{1}} \cdots S^{\varrho_{k}} S^{\varrho} | p_{\varrho}).
\end{equation}

Die einzige Spezialisierung, die in diesen Formeln enthalten ist, daß eine Reihe $x$ eintritt, läßt sich nicht umgehen; bei ganz allgemeinen Symbol- und Variablenreihen gelangen wir zu einer ``Zerlegungsidentität,'' einer Identität, bei der eine Reihe $p_{\varrho}$ in die einzelnen Reihen $x$ zerlegt wird. Den einfachsten Spezialfall der Zerlegungsidentität benutzte schon Clebsch in seiner ``Fundamentalaufgabe der Invariantentheorie'' zur Ableitung der Elementarkovarianten bei den Reihenentwicklungen.

Mit Hilfe dieser Identitäten, die den Übergang von den nicht expliziten Determinantenausdrücken zum Matrizenprodukt vermitteln, läßt sich nun in der Tat nachweisen, daß mit der Matrizendarstellung die gewünschte explizite symbolische Darstellung gewonnen ist.

Für den Faltungsprozeß aber gilt ein ganz analoger Satz wie im ternären Gebiet, auf den sich auch analog die Reduktionssätze aufbauen lassen. Bezeichnet man in (1) die Zahl $\lambda$ als \textit{Defekt} der Faltung, und Faltungen, für die $\varrho = \tau$, als \textit{kogredient}, so lassen sich alle Faltungen erzeugen durch sukzessive Anwendung der kogredienten Faltungen vom Defekt~$1$. Der Beweis dieses Satzes beruht wesentlich darauf, daß man die allgemeinsten Formen durch ``Normalformen'' ersetzen kann, d.h.\ durch Formen, die bei Anwendung aller invarianten Differentialoperationen identisch verschwinden. Im quaternären Gebiet hat diese letztere Tatsache Herr Mertens\footnote{Über invariante Gebilde quaternärer Formen. \textit{Wiener Berichte.} Bd.~98. 1889.} bewiesen; unsere Kenntnis der allgemeinsten Differentialoperationen --- die bekanntlich aus den Faltungsprozessen durch Ersetzen der Symbolreihen durch Differentialsymbole hervorgehen --- gestattet uns, unter Anwendung der Zerlegungsidentität den Mertenschen Beweis fast wörtlich auf das $n$-äre Gebiet zu übertragen.

\newpage
% ============================================
% PAPER 4: Pages 118--150
% ============================================

\begin{center}
\textbf{4. Zur Invariantentheorie der Formen von $n$ Variabeln}

\smallskip
\textit{Journal f.~d.~reine u.~angew.\ Math.\ 139 (1911), S.~118--154}
\end{center}

\section*{Einleitung.}
\addcontentsline{toc}{section}{Einleitung}

In der projektiven Invariantentheorie der Formen von $n$ Variabeln sind die an das Hauptproblem --- den Endlichkeitsnachweis des Formensystems --- sich anschließenden Fragen kaum behandelt, wenn man über das binäre und ternäre Gebiet hinausgeht. Es sind dies die Fragen nach der gegenseitigen Abhängigkeit der Formen und nach den Methoden zur Beherrschung dieses Formenzusammenhangs. Diese Methoden beruhen wesentlich --- wie dies im binären und ternären Gebiet vollständig nachgewiesen ist --- auf zwei von Herrn Study als ``Fundamentalsätze der symbolischen Methode'' bezeichneten Sätzen, nämlich dem Satz von der symbolischen Darstellbarkeit der Formen und dem Satz von den symbolischen Identitäten\footnote{E.\ Study, \textit{Methoden zur Theorie der ternären Formen} (Leipzig, Teubner, 1889). II. \S\ 6. --- Über das Auftreten dieser Fundamentalsätze auch in der Invariantentheorie spezieller Transformationsgruppen vgl.\ Study, \textit{Das Apollonische Problem} (Math.\ Ann.\ Bd.~49 [1897]) und \textit{Zur Differentialgeometrie der analytischen Kurven} (Transactions of the American Math.\ Society, Bd.~1 [1909]).}. Der letztere Satz sagt aus, daß man alle zwischen ganzen invarianten Bildungen bestehenden Relationen erhält durch sukzessive Anwendung einer kleinen Anzahl von symbolischen Identitäten; daß also die symbolische Methode, und nur diese, die Kenntnis des Formenzusammenhangs gibt, ohne aus dem Gebiet des Invarianten herauszutreten.

Auf den Grund der Schwierigkeit beim Übergang zu $n$ Variabeln hat schon Herr Gordan in seinem Erlanger Programm\footnote{P.\ Gordan, \textit{über das Formensystem binärer Formen.} S.~46. (Leipzig, Teubner, 1875).} hingewiesen; er liegt darin, daß der erste Satz der symbolischen Methode in seiner allgemeinen Fassung nicht mehr gilt. Es treten im $n$-nären Gebiet bekanntlich Variabelnreihen höherer Stufe $p_{\varrho}$\textsuperscript{,}\footnote{Für die Bezeichnung vgl.\ \S\ 1.} auf und analog Symbolreihen höherer Stufe $S^{\varrho}$; und zwar gelangt man zu diesen Reihen sowohl, wenn man, wie ursprünglich Clebsch\footnote{A.\ Clebsch, \textit{über eine Fundamentalaufgabe der Invariantentheorie.} (Abh.\ der Ges.\ d.\ Wiss.\ zu Göttingen, Bd.~XVII, 1872).}, ausgeht von Formen mit beliebig vielen Reihen $p_{\varrho}$, als auch wenn man ausgeht von Formen, die nur beliebig viele kogrediente Reihen $x$ enthalten, nach dem Vorgang von F.\ Deruyts und A.\ Capelli\footnote{Vgl.\ A.\ Capelli, \textit{Lezioni sulla teoria delle forme algebriche} (Napoli 1902), \S\ 22--24, und die dort angeführte Literatur.}. Eine symbolische Darstellung, die die Reihen $p_{\varrho}$, $S^{\varrho}$ explizit\textsuperscript{,}\footnote{Genauere Definition der ``expliziten Darstellung'' in \S\ 2.} enthielte, existiert nun nicht; es ist nötig, die $p_{\varrho}$, $S^{\varrho}$ in die einzelnen Reihen $x$ und die dazu kontragredienten $s$ aufzulösen und diese in verschiedene Determinanten zu verteilen. Nun läßt sich allerdings mittels Differentialbedingungen (vgl.\ Formel (19.)) verifizieren, daß ein gegebenes Determinantenaggregat die Eigenschaft besitzt, nur von den Reihen $p_{\varrho}$, $S^{\varrho}$ abzuhängen; ein Überblick aber über die Gesamtheit der invarianten Bildungen und über die Prozesse zu ihrer Erzeugung läßt sich auf diese Art nicht gewinnen\textsuperscript{,}\footnote{Auch eine von Herrn Study herrührende rechnerisch sehr wertvolle Weiterbildung der Symbolik (mitgeteilt von F.\ Engel, Ber.\ d.\ K.\ Sächs.\ Ges.\ d.\ Wiss., math.-phys.\ Kl., 1900, S.~126, und für das quaternäre Gebiet von E.\ Study, \textit{Geometrie der Dynamen}, S.~135) ist keine explizite Darstellung in unserm Sinne. Sie dürfte z.~B.\ schwerlich zu den unsymbolischen Differentiationsprozessen zur Erzeugung der Formen führen.}.

Wir zeigen nun in \S\S\ 2 und 6, wie man unter Anwendung des ``Satzes von den korrespondierenden Matrizen''\footnote{Abgesehen von Graßmanns Ausdehnungslehre von 1862, Nr.~112 zuerst bei A.\ Brill, \textit{über zwei Berührungsprobleme,} \S\ 2 (Math.\ Ann.\ Bd.~4. 1871).} den ersten Fundamentalsatz in seiner Allgemeinheit wieder erhält, indem man die Darstellung durch Determinanten ersetzt durch eine Darstellung durch ``Matrizenprodukte''. In \S\ 2 wird gezeigt, daß ein jedes die Reihen $S^{\varrho}$,$p_{\varrho}$ explizit enthaltende Matrizenprodukt eine invariante Bildung der Grundform ist; die in \S\ 6 gegebene Umkehrung, daß sich jede invariante Bildung explizit durch Matrizenprodukte darstellen läßt, gelingt erst durch Übertragung des zweiten Fundamentalsatzes (\S\S\ 3--5).

Dieser Satz ist bekannt für Formen, die nur Variablenreihen $x$ enthalten\footnote{E.\ Pascal, \textit{Giorn.\ di mat.}\ 26 (1888), S.~33, 102; \textit{Rendic.\ d.\ Accad.\ d.\ Line.} (4) 4 (1888), S.~119; ib.\ Mem.\ (4) 5 (1888), S.~375.}. Das allgemeine Problem, die Gesamtheit der unzerlegbaren Identitäten aufzustellen, bei denen alle Reihen $S^{\varrho}$,$p_{\varrho}$ als unzerlegbar betrachtet werden, ist von Herrn Study formuliert\footnote{``Methoden \ldots'', S.~204, Anmerkung 17).}; er sieht zugleich in der Anzahl der auftretenden Identitäten ein Maß für die Schwierigkeit der Methoden im $n$-nären Gebiet. Indem es nun gelingt, die Gesamtheit der Identitäten in eine einzige Formel (36.) zusammenzufassen, ist diese Schwierigkeit wenigstens auf ein Mindestmaß herabgedrückt. Bei der Anwendung werden allerdings häufig gewisse ausgezeichnete Spezialfälle von (36.) auftreten, wie (20.), (21.), die dadurch charakterisiert sind, daß sie die explizite Darstellung zulassen, ohne Substitution neuer Reihen; oder Formel (31.), die als ``Zerlegungsidentität'' bezeichnet wird, da eine Reihe $p_{\varrho}$ scheinbar in die Reihen $y$ zerlegt auftritt.

Auf den beiden Fundamentalsätzen läßt sich nun leicht die weitere Theorie aufbauen. Der erste Satz gibt direkt die Prozesse zur Erzeugung der Formen, den symbolischen Faltungsprozeß und die unsymbolischen Differentiationsprozesse\textsuperscript{,}\footnote{In der Arbeit \textit{``Invariante Gebilde von Nullsystemen''} (Sitzungsber.\ d.\ Ak.\ d.\ Wiss.\ Wien, Bd.~97, Abt.\ IIa, 1888) ist Herr Mertens auf anderem, nicht direkt verallgemeinerungsfähigem Wege zum quaternären Faltungsprozeß gelangt. Die dort auftretende alternierende Invariante von 6 Reihen $S^{2}$ unterscheidet sich von der bei uns bei einmaliger Anwendung der Substitution (11.) entstehenden durch ein Aggregat gerader Invarianten. Für einen Spezialfall findet sich der quaternäre Faltungsprozeß auch bei Herrn P.\ Gordan, \textit{Das simultane System von zwei quadratischen quaternären Formen} (Math.\ Ann.\ Bd.~56. 1903).}, während die Zerlegungsidentität unter diesen Prozessen alle äquivalenten ausschaltet (\S\ 6). Die Kenntnis der Differentiationsprozesse ergibt weiter die Übertragung des Begriffes der ``Normalform'' auf das $n$-näre Gebiet, und mittels eines im quaternären Gebiet von Herrn Mertens\footnote{F.\ Mertens, \textit{Über invariante Gebilde quaternärer Formen} (siehe Einleitung), zitiert als ``Mertens''.} angegebenen Beweisverfahrens den Satz, daß sich ein System von invarianten Bildungen ersetzen läßt durch ein äquivalentes System von Normalformen (\S\ 7). Unter dieser letzteren Voraussetzung habe ich in meiner Dissertation\footnote{\textit{Über die Bildung des Formensystems der ternären biquadratischen Form.} (Dieses Journal, Bd.~134. 1908.)} gezeigt, daß sich der ternäre Faltungsprozeß erzeugen läßt durch sukzessive Anwendung von zwei Grundfaltungen. Auf Grund dieses Satzes und der Theorie der Reihenentwicklung habe ich dann unter Einführung des Begriffes der ``Formenreihe'' die hauptsächlichsten Reduktionssätze für Formensysteme entwickelt. Ganz analog läßt sich im $n$-nären Gebiet der Faltungsprozeß aus $n-1$ Grundfaltungen erzeugen (\S\ 8) und lassen sich die Reduktionssätze aufstellen (\S\ 9).

\subsection*{Nachbemerkung.}
\addcontentsline{toc}{subsection}{Nachbemerkung}

Gelegentlich der Salzburger Mitteilung wurde ich durch Herrn Prof.\ E.\ Müller, Wien, darauf aufmerksam gemacht, daß das dieser Arbeit zugrundeliegende Rechnen mit Matrizenprodukten im Prinzip identisch ist mit dem Kalkül von Graßmanns Ausdehnungslehre\footnote{Hermann Graßmanns gesammelte mathematische und physikalische Werke. Ersten Bandes zweiter Teil: \textit{Die Ausdehnungslehre von 1862.} In Gemeinschaft mit H.\ Graßmann d.\ Jüngeren herausgegeben v.\ Fr.\ Engel (Leipzig, Teubner 1896).}. In der Tat läßt sich Graßmanns ``kombinatorisches Produkt'' $[S^{\varrho_{1}} S^{\varrho_{2}} \cdots S^{\varrho_{k}}]$ auffassen als eine durch diese Reihen gegebene Matrix (vgl.\ \S\ 2), und zwar das ``progressive Produkt'' als die Matrix der Reihen selbst, das ``regressive'' als die Matrix der zugeordneten Reihen $S^{n-\varrho_{1}}$, $S^{n-\varrho_{2}}$, \ldots $S^{n-\varrho_{k}}$, während die dem ``gemischten Produkt'' entsprechende Matrix entsteht, indem man mehrere Reihen $S$ durch die Substitution (9.) zu neuen Reihen $P$,$Q$ zusammenfaßt. Unsere ``zugeordnete Reihe'' entspricht dabei Graßmanns ``Ergänzung'', unser ``Matrizenprodukt'' einem ``gemischten Produkt von der Stufenzahl~$0$''.

Nach dieser Zuordnung wird die in der Ausdehnungslehre 1862, Nr.~113, gegebene Entwicklung identisch mit einer zu unserer Formel (32.) dualistischen Formel. Im Spezialfall $\varrho=1$ geht (32.) in (17.), die dualistische Formel in (16.) über. Aus diesen Spezialfällen der Graßmannschen Entwicklung hat neuerdings Herr Müller\footnote{E.\ Müller, \textit{Beiträge zur Graßmannschen Ausdehnungslehre.} 1.\ Mitteilung: Einige allgemeine Sätze. (Sitzungsber.\ d.\ Ak.\ d.\ Wiss.\ Wien; Bd.~118. Abt.\ IIa. Juli 1909), Satz 16 und 18.} die oben erwähnten Identitäten (20.), (21.) abgeleitet.

Unsere Ableitung gibt im Unterschied von der Graßmann-Müllerschen zugleich den Nachweis der Vollständigkeit der Identitäten, was für die hier in Betracht kommenden Fragen als das Wesentlichste erscheint, und sie geht von (20.), (21.) weiter zu den allgemeinsten unzerlegbaren Identitäten (36.), die bis jetzt noch nicht bemerkt zu sein scheinen.

\section*{\S\ 1. Bezeichnungen und Definitionen. Zusammenfassung bekannter Resultate.}
\addcontentsline{toc}{section}{\S\ 1. Bezeichnungen und Definitionen}

Wir bezeichnen, wie üblich, als Variablenreihe $x$ die Reihe der Elemente $x_{1}, x_{2}, \ldots x_{n}$, und analog als Variablenreihe $p_{\varrho}$ ($\varrho = 1, 2, \ldots n-1$) die Reihe der $\binom{n}{\varrho}$ Elemente $p_{i_{1} i_{2} \cdots i_{\varrho}}$, wobei $p_{i_{1} i_{2} \cdots i_{\varrho}}$ die aus der Matrix der $\varrho$ Reihen $x^{(1)}, x^{(2)}, \ldots x^{(\varrho)}$ gebildeten $\binom{n}{\varrho}$ Determinanten
\begin{equation}\tag{1}
x^{(1)}_{\alpha_{1}} x^{(1)}_{\alpha_{2}} \cdots x^{(1)}_{\alpha_{\varrho}} \qquad (1 \leq \alpha_{1} < \alpha_{2} < \cdots < \alpha_{\varrho} \leq n)
\end{equation}
bedeuten. Nach dem Satze über die korrespondierenden Matrizen\footnote{Vgl.\ etwa: Gordan-Kerschensteiner, \textit{Vorlesungen über Invariantentheorie,} Bd.~II, S.~99.} ist einer jeden Reihe $p_{\varrho}$ eine zweite $q_{n-\varrho}$ ein-eindeutig zugeordnet durch die für jedes Element der Reihe geltende Beziehung:
\begin{equation}\tag{1$'$}
p_{\varrho} \cdot q_{n-\varrho} = \begin{vmatrix} x^{(1)}_{1} & \cdots & x^{(1)}_{n} \\ \vdots & & \vdots \\ x^{(\varrho)}_{1} & \cdots & x^{(\varrho)}_{n} \\ u^{(1)}_{1} & \cdots & u^{(1)}_{n} \\ \vdots & & \vdots \\ u^{(n-\varrho)}_{1} & \cdots & u^{(n-\varrho)}_{n} \end{vmatrix},
\end{equation}
worin die $\alpha$ und $\beta$ unter sich gut geordnet sind und zusammen eine volle Permutation der Zahlen $1$ bis $n$ ergeben. Die Reihe $q_{n-\varrho}$ besteht aus den $\binom{n}{\varrho}$ Determinanten $(n-\varrho)$-ten Grades, die gebildet sind aus der Matrix von $n-\varrho$ den $x$ kontragredienten Reihen $u^{(1)}, u^{(2)}, \ldots u^{(n-\varrho)}$. Diese Reihen erfüllen die $\varrho(n-\varrho)$ Bedingungsgleichungen:
\begin{equation}\tag{2}
(u^{(k)} | x^{(l)}) = \sum_{\alpha=1}^{n} u^{(k)}_{\alpha} x^{(l)}_{\alpha} = 0, \qquad (k = 1, 2, \ldots n-\varrho; \, l = 1, 2, \ldots \varrho).
\end{equation}

Die allgemeinsten Formen lassen sich nach Clebsch\footnote{Clebsch, a.\ a.\ O. \S\ 4.} (und ebenso nach den späteren Untersuchungen von F.\ Deruyts und A.\ Capelli, vgl.\ Einleitung) zurückführen auf solche, die von jeder der $n-1$ Reihen $p_{\varrho}$ höchstens eine enthalten, die sich also symbolisch folgendermaßen darstellen lassen:
\begin{equation}\tag{3}
f = (S^{\varrho_{1}} | p_{\varrho_{1}})^{m_{1}} (S^{\varrho_{2}} | p_{\varrho_{2}})^{m_{2}} \cdots (S^{\varrho_{k}} | p_{\varrho_{k}})^{m_{k}},
\end{equation}
wobei wir in Analogie mit (2.) gesetzt haben:
\begin{equation*}
(S^{\varrho} | p_{\varrho}) = \sum_{i_{1} < \cdots < i_{\varrho}} S^{\varrho}_{i_{1} \cdots i_{\varrho}} p_{i_{1} \cdots i_{\varrho}}.
\end{equation*}

Die Symbolreihen $S^{\varrho}$ lassen sich infolge der zwischen den Elementen einer Reihe $p_{\varrho}$ bestehenden nichtlinearen identischen Beziehungen so normieren, daß sie ebendenselben Identitäten genügen, daß sie sich also verhalten wie die Determinanten einer $\varrho$-reihigen Matrix, deren Reihen $s^{(1)}, s^{(2)}, \ldots s^{(\varrho)}$ zu den $x$ kontragredient sind\footnote{Clebsch, a.\ a.\ O. \S\ 5.}. Es läßt sich also einer jeden Reihe $S^{\varrho}$ eine andere $S^{n-\varrho}$ ein-eindeutig zuordnen durch die Beziehung:
\begin{equation}\tag{4}
S^{\varrho}_{\alpha_{1} \cdots \alpha_{\varrho}} = (-1)^{\sum \alpha + \sum \beta} S^{n-\varrho}_{\beta_{1} \cdots \beta_{n-\varrho}},
\end{equation}
wo wieder die $\alpha$ und $\beta$ unter sich gut geordnet sind und zusammen die Zahlen $1$ bis $n$ gerade erschöpfen. Die Elemente $S^{n-\varrho}$ der Reihe $S^{n-\varrho}$ verhalten sich wie die aus $(n-\varrho)$ --- den $x$ kogredienten --- Reihen $\sigma^{(1)}, \sigma^{(2)}, \ldots \sigma^{(n-\varrho)}$ gebildeten Determinanten, und die Reihen $s$ und $\sigma$ sind durch die $\varrho(n-\varrho)$ Bedingungen verknüpft:
\begin{equation}\tag{5}
(s^{(k)} | \sigma^{(l)}) = 0, \qquad (k = 1, 2, \ldots \varrho; \, l = 1, 2, \ldots n-\varrho).
\end{equation}

Infolge der Beziehungen (1.) und (4.) ist ein jeder Faktor von (3.) der äquivalenten vierfachen symbolischen Darstellung fähig:
\begin{equation*}
(S^{\varrho} | p_{\varrho}) = (S^{n-\varrho} | q_{n-\varrho}) = (S^{\varrho} | q_{n-\varrho}) = (S^{n-\varrho} | p_{\varrho}),
\end{equation*}
wobei, wie stets im folgenden, $(S^{\varrho} q_{n-\varrho})$ und $(S^{n-\varrho} p_{\varrho})$ abkürzend für die Determinanten $(s^{(1)} \cdots s^{(\varrho)} u^{(1)} \cdots u^{(n-\varrho)})$ und $(\sigma^{(1)} \cdots \sigma^{(n-\varrho)} x^{(1)} \cdots x^{(\varrho)})$ gesetzt sind; durch Laplacesche Entwicklung erweisen sich diese als nur abhängig von den $p_{\varrho}$.

Die einer Symbol- (bzw.\ Variabeln-) reihe zukommende charakteristische Zahl $\varrho(n-\varrho)$ bezeichnen wir nach Clebsch als \textit{Gewicht} der Reihe\footnote{Vgl.\ einen Nachtrag von Study in Graßmanns Werken, ersten Bandes zweiter Teil, S.~510 (Bedingung, daß eine Größe $S^{\varrho}$ einfach ist).}, die Zahl $\varrho$, die die Anzahl der Reihen $s$ angibt, als \textit{oberen Gewichtsindex}, die Zahl $n-\varrho$, die die Anzahl der Reihen $\sigma$ angibt, als \textit{unteren Gewichtsindex}. Aus den Beziehungen (4.) folgt, daß eine Symbolreihe in ein und derselben Relation einmal mit oberem, einmal mit unterem Gewichtsindex auftreten kann; dasselbe gilt für das Auftreten der zusammengehörigen Reihen $p_{\varrho}$ und $q_{n-\varrho}$.

Es sei noch bemerkt, daß wir allgemein bezeichnen:\\
Variabelnreihen höherer Stufe, deren zugehörige Matrix aus Reihen $x$ besteht, durch $p$; solche, deren Matrix aus Reihen $u$ besteht, durch $q$;\\
die den $x$ kogredienten Symbolreihen durch kleine griechische Buchstaben: $\varrho$, $\tau$, $\sigma$, $\beta$;\\
die den $u$ kogredienten Symbolreihen durch kleine lateinische Buchstaben: $S$, $T$, $R$, \ldots;\\
alle übrigen Symbolreihen durch große lateinische Buchstaben mit Gewichtsindex: $S^{\varrho}$, $T^{\tau}$, $R^{\varrho}$, \ldots oder $P_{\varrho}$, $Q_{\varrho}$, \ldots.

Entsprechend oberem oder unterem Gewichtsindex schreiben wir auch die einzelnen Elemente einer Reihe mit oberen oder unteren Indizes, wie z.B.\ in (4.).

\section*{\S\ 2. Darstellung durch Matrizenprodukte.}
\addcontentsline{toc}{section}{\S\ 2. Darstellung durch Matrizenprodukte}

Wir verstehen im folgenden stets unter ``Matrizenprodukt'' die Summe über die Produkte entsprechender Determinanten zweier Matrizen von gleicher Reihenzahl, wo die Reihen der zweiten Matrix denjenigen der ersten kontragredient sind. Dabei können die Reihen jeder Matrix: 1.\ einzeln gegeben sein, 2.\ wie in (6.) zu einer Reihe $S^{\varrho}$ zusammengefaßt, 3.\ zu verschiedenen Reihen $S^{\varrho}$, $S^{\tau}$, \ldots\ vereinigt sein. Bezeichnet soll das Matrizenprodukt durch das in (6.) angewandte Symbol $(\quad | \quad)$ werden. Die Darstellung
\begin{equation}\tag{6}
f = (S^{\varrho} | p_{\varrho}) = (S^{n-\varrho} | q_{n-\varrho})
\end{equation}
ist eine spezielle Anwendung einer zweiten Fassung des Satzes über die korrespondierenden Matrizen: ``Einem jeden $\varrho$-reihigen Matrizenprodukt $(v^{(1)} \cdots v^{(\varrho)} | y^{(1)} \cdots y^{(\varrho)})$ ist eine Determinante zugeordnet; und umgekehrt läßt sich jeder Determinante ein Matrizenprodukt von vorgegebener Reihenanzahl $\varrho$ ($\varrho = 1, 2, \ldots n-1$) zuordnen''\footnote{Für $\varrho > n$ verschwindet bekanntlich das Matrizenprodukt; für $\varrho = n$ geht es in das Produkt zweier Determinanten über.}. In der Tat hat man nur die Reihen $y$ zu einer Reihe $p_{\varrho}$ zusammenzufassen (bzw.\ die Reihen $v$ zu einer Reihe $q_{\varrho}$) und die durch (1.) gegebenen Substitutionen auszuführen, um bei gegebenem Matrizenprodukt zur Determinante zu gelangen und umgekehrt. Der Wert der Determinante ist dabei allerdings nicht durch ihre einzelnen Elemente gegeben, sondern ergibt sich erst durch Laplacesche Entwicklung. Da sich also Determinante und Matrizenprodukt als äquivalent erweisen, drückt sich der Satz von der symbolischen Darstellbarkeit der Formen (erster Fundamentalsatz) auch so aus:

\medskip
\textbf{Satz I:} \textit{``Alle invarianten Bildungen lassen sich symbolisch durch Matrizenprodukte darstellen, und umgekehrt sind alle Matrizenprodukte invariante Bildungen''}\footnote{Invariante Bildungen vorgegebener Grundformen sind natürlich nur solche Aggregate von Matrizenprodukten, die von den Symbolreihen $S^{\varrho}$, \ldots\ abhängen.}.
\medskip

Die Darstellung durch Matrizenprodukte erlaubt, die durch die Determinantendarstellung bedingte wesentliche Schwierigkeit der nicht expliziten Darstellung zu überwinden\footnote{Vgl.\ die Einleitung.}. Wir verstehen dabei unter ``expliziter Darstellung'' eine solche, in die nur die Reihen $S^{\varrho}$, $T^{\tau}$, \ldots\ selbst eingehen, oder eindeutig aus diesen ableitbare Reihen $R^{\varrho}$, in der aber die einzelnen Teilreihen $s$, $t$, \ldots\ nicht mehr auftreten. Wir werden in \S\ 6 für alle invarianten Bildungen, die nur von den Symbolreihen $S^{\varrho}$, \ldots\ $T^{\tau}$ abhängen, eine in diesen Symbolreihen explizite Darstellung erhalten, und damit die Übertragung der erzeugenden Prozesse, Faltungsprozeß und Differentiationsprozesse, auf das $n$-näre Gebiet. Die Möglichkeit der expliziten Darstellung erhellt daraus, daß nur die einfachsten Matrizenprodukte gerade den bei der Determinantendarstellung auftretenden Determinanten zugeordnet werden können, daß also alle übrigen Matrizenprodukte Zusammenfassungen verschiedener Determinanten zu einem einzigen Ausdruck sind. Denn löst man die Symbol- und Variabelnreihen in die einzelnen Reihen $s$,$u$,$x$ auf, so muß nach dem ersten Fundamentalsatz in seiner gewöhnlichen Form notwendig ein Aggregat von Determinanten entstehen. Der Beweis dafür, daß umgekehrt alle Determinantenaggregate, die invariante Bildungen vorgegebener Grundformen darstellen, sich zu expliziten Matrizenprodukten zusammenfassen lassen, ergibt sich unter Benutzung des zweiten Fundamentalsatzes (\S\ 6).

Um nun aus zwei Symbolreihen $S^{\varrho}$, $T^{\tau}$ unter Zuhilfenahme von Variablenreihen eine alle Reihen explizit enthaltende invariante Bildung der Grundform $(S^{\varrho} | p_{\varrho})(T^{\tau} | p_{\tau})$ zu erhalten (Faltungsprozeß), haben wir nur ein Matrizenprodukt dritter Art nach der am Anfang des Paragraphen gegebenen Definition zu bilden:
\begin{equation}\tag{7}
\Phi = (G_{n-\varrho-\lambda} S^{\varrho} T^{\tau} | p_{\varrho+\lambda} p_{\tau-\lambda}) \qquad (\lambda = 1, 2, \ldots \tau, n-\varrho)\footnote{Die Zahl $\lambda$ läuft stets bis zu der kleineren der beiden zuletzt angegebenen Zahlen.}
\end{equation}

Entwickelt gibt (7.):
\begin{equation}\tag{8}
\Phi = \sum \begin{vmatrix} s^{(1)}_{\alpha_{1}} & \cdots & s^{(1)}_{\alpha_{\varrho}} & \sigma^{(1)}_{\beta_{1}} & \cdots & \sigma^{(1)}_{\beta_{\lambda}} \\ \vdots & & \vdots & \vdots & & \vdots \\ s^{(\varrho)}_{\alpha_{1}} & \cdots & s^{(\varrho)}_{\alpha_{\varrho}} & \sigma^{(\varrho)}_{\beta_{1}} & \cdots & \sigma^{(\varrho)}_{\beta_{\lambda}} \\ t^{(1)}_{\alpha_{1}} & \cdots & t^{(1)}_{\alpha_{\varrho}} & \tau^{(1)}_{\beta_{1}} & \cdots & \tau^{(1)}_{\beta_{\lambda}} \\ \vdots & & \vdots & \vdots & & \vdots \\ t^{(\tau)}_{\alpha_{1}} & \cdots & t^{(\tau)}_{\alpha_{\varrho}} & \tau^{(\tau)}_{\beta_{1}} & \cdots & \tau^{(\tau)}_{\beta_{\lambda}} \end{vmatrix} \cdot \begin{vmatrix} x^{(1)}_{\alpha_{1}} & \cdots & x^{(1)}_{\alpha_{\varrho}} & u^{(1)}_{\beta_{1}} & \cdots & u^{(1)}_{\beta_{\lambda}} \\ \vdots & & \vdots & \vdots & & \vdots \\ x^{(\varrho)}_{\alpha_{1}} & \cdots & x^{(\varrho)}_{\alpha_{\varrho}} & u^{(\varrho)}_{\beta_{1}} & \cdots & u^{(\varrho)}_{\beta_{\lambda}} \\ x^{(\varrho+1)}_{\alpha_{1}} & \cdots & x^{(\varrho+1)}_{\alpha_{\varrho}} & u^{(\varrho+1)}_{\beta_{1}} & \cdots & u^{(\varrho+1)}_{\beta_{\lambda}} \\ \vdots & & \vdots & \vdots & & \vdots \\ x^{(\varrho+\lambda)}_{\alpha_{1}} & \cdots & x^{(\varrho+\lambda)}_{\alpha_{\varrho}} & u^{(\varrho+\lambda)}_{\beta_{1}} & \cdots & u^{(\varrho+\lambda)}_{\beta_{\lambda}} \end{vmatrix}
\end{equation}
d.h.\ einen nur von den Reihen $S$,$T$,$q$,$p$ abhängigen Ausdruck. Es ist dabei $\varepsilon = \pm 1$\footnote{Diese Bezeichnung soll durchweg beibehalten werden.} und die Summe so zu verstehen, daß die Indizes einer jeden Reihe $S$, \ldots\ unter sich gut geordnet sind, und für jede Kombination $\alpha_{1} \cdots \alpha_{\varrho}$, $\beta_{1} \cdots \beta_{\lambda}$ die $\alpha$ sowohl wie die $\beta$ einzeln alle dann noch möglichen Permutationen der Zahlen durchlaufen.

Die in (7.) eingehende Zahl $\lambda$ bezeichnen wir als \textit{Defekt} der Faltung im Falle $\varrho > \tau$; der Grund dieser letzteren Einschränkung ergibt sich in \S\ 6.

Die Determinanten der beiden in (7.) eingehenden Matrizen lassen sich als Elemente neuer Reihen $Q^{\varrho+\lambda}$,$P_{\tau-\lambda}$ betrachten, deren zugeordnete $Q_{\varrho}$, $P^{\tau}$ seien. Diese Reihen $Q$,$P$ lassen sich nach (8.) durch die ursprünglichen $S$ und $q$, bzw.\ $T$ und $p$ ausdrücken. Wir deuten dies an durch die Gleichungen:
\begin{equation}\tag{9}
Q^{\varrho} = [q_{\varrho} \cdot S^{\varrho}]_{\varrho+\lambda}, \qquad P^{\tau} = [p_{\tau} \cdot T^{\tau}]_{\tau-\lambda}
\end{equation}
oder auch
\begin{equation*}
q_{\varrho} = [q_{\varrho+\lambda} \cdot S^{\varrho}]_{\varrho}, \qquad p_{\tau} = [p_{\tau-\lambda} \cdot T^{\tau}]_{\tau}.
\end{equation*}

Unter Berücksichtigung von (4.) gilt dann, analog (6.), für das Matrizenprodukt (7.) die äquivalente vierfache symbolische Darstellung:
\begin{equation}\tag{10}
(G_{n-\varrho-\lambda} S^{\varrho} T^{\tau} | p_{\varrho+\lambda} p_{\tau-\lambda}) = (G_{\varrho} S^{\varrho} | p_{\varrho} P^{\tau}) = (G_{n-\varrho} S^{\varrho} | q_{n-\varrho} Q^{\varrho+\lambda}) = \cdots
\end{equation}

Eine weitere Substitution neuer Reihen läßt sich an (7.) vornehmen, indem man setzt:
\begin{equation}\tag{11}
(G_{n-\varrho-\lambda} S^{\varrho} T^{\tau} | p_{\varrho+\lambda} p_{\tau-\lambda}) = (R^{n} | p_{\varrho} p_{\tau}).
\end{equation}
Auch hier zeigt (8.), daß die Produkte der Elemente der Reihen $R$ sich eindeutig durch die Reihen $S$ und $T$ ausdrücken lassen. Die Substitutionen (9.) und (11.) werden insbesondere dann anzuwenden sein, wenn mit (7.) eine weitere Faltung mit anderen Reihen vorgenommen wird.

Die Matrizendarstellung setzt uns instand, die quadratischen Relationen zwischen den Elementen einer Reihe $p_{\varrho}$ (aus denen bekanntlich alle übrigen folgen), also nach \S\ 1 die in den Koeffizienten der Grundform linearen Beziehungen, denen die Reihen $S^{\varrho}$ genügen, invariant aufzustellen. Es sind die Identitäten:
\begin{equation}\tag{12}
(G_{n-\varrho-\lambda} S^{\varrho}_{1} \cdots S^{\varrho}_{\lambda+1} | p_{\varrho} p_{\varrho+1} \cdots p_{\varrho+\lambda}) = 0, \qquad (\lambda = 1, 2, \ldots \varrho-1)
\end{equation}
wo die $p$ und $q$ als willkürliche Größen zu betrachten sind.

Wir werden in \S\ 5 sehen, daß die Identitäten mit ungerader Defektzahl $\lambda$ eine Folge derjenigen von höherem Defekt sind. Die Beziehungen (12.) sind eine direkte Folge der Relationen (5.); denn man hat:
\begin{equation*}
(G_{n-\varrho} S^{\varrho}_{i} | q_{n-\varrho} Q^{n}) = 0 \quad \text{für jedes } i,
\end{equation*}
und die Anwendung des Produktsatzes gibt infolge von (5.) eine verschwindende $(n-\lambda)$-reihige Determinante.

Die Identitäten (12.) sagen aus, daß es zur Aufsuchung aller Typen von Faltungen genügt, das Produkt von in jeder Variabelnreihen linearen Formen zu falten; oder auch, nach \S\ 6, daß die normierte Form $(S^{\varrho} | p_{\varrho})^{m}$ keine in den Koeffizienten lineare invariante Bildung besitzt.

Die Bedingungen (12.) sind auch hinreichend zur Charakterisierung der Reihen $p_{\varrho}$. Da nämlich die Eigenschaft, daß die einzelnen Elemente von $p_{\varrho}$ Determinanten einer Matrix sind, eine invariante ist, muß sie sich auch invariant ausdrücken lassen; die Bedingungen (12.) stellen aber nach Satz VI die größtmögliche invariante Einschränkung für die $p_{\varrho}$ dar, nämlich die, daß die aus der Reihe $p_{\varrho}$ als Koeffizienten gebildete Linearform überhaupt keine invarianten Bildungen besitzt.


\section*{\S\ 3. Die symbolischen Identitäten.}
\addcontentsline{toc}{section}{\S\ 3. Die symbolischen Identitäten}

Wir haben nun den zweiten Fundamentalsatz der symbolischen Methode zu übertragen, nämlich die Gesamtheit der irreduziblen unzerlegbaren Identitäten aufzustellen, bei denen alle Reihen $S^{\varrho}$, $p_{\varrho}$ als unzerlegbar angesehen werden. Wir setzen dabei noch folgendes fest: Entsprechend der Bedeutung von Symbol- und Variabelnreihen sollen solche Identitäten als zusammengesetzt betrachtet werden, die entstehen, wenn man die Variablenreihe $p_{\varrho}$ zu $p_{\varrho'} p_{\varrho''}$ spezialisiert und auf die entstehenden Ausdrücke eine der Identitäten des Systems anwendet; während andrerseits Identitäten, die $k$ Symbolreihen enthalten, als unzerlegbar gelten sollen, auch wenn sie durch den obigen Prozeß aus Identitäten mit weniger Symbolreihen hervorgehen. Die Reihen $S^{\varrho}$,$p_{\varrho}$ brauchen nicht notwendig explizit in die Identitäten einzugehen; zu ihrer Auffassung als unzerlegbare Größen genügt der Nachweis, daß die auftretenden Matrizenprodukte nur von diesen Reihen abhängen. Wir werden aber zeigen, daß in diesem Falle, teilweise unter Anwendung der Substitution (11.), sich stets eine explizite Darstellung erzielen läßt. Die allgemeinen Identitäten erhalten wir mittels des Prinzipes der Matrizendarstellung aus den von Herrn Pascal gegebenen speziellen, die nur Reihen $x$ und $u$ enthalten und eine direkte Übertragung der von Herrn Study für das ternäre Gebiet gegebenen darstellen:
\begin{align}\tag{13}
\sum_{i=1}^{n} (a^{(1)} \cdots a^{(\varrho)} | x^{(i)}) (a^{(\varrho+1)} \cdots a^{(k)} | x^{(i)}) &= (a^{(1)} \cdots a^{(k)} | x), \\
\tag{14}
(a^{(1)} | x) (a^{(2)} | x) \cdots (a^{(n)} | x) &= (a^{(1)} \cdots a^{(n)})(x^{(1)} \cdots x^{(n)}), \\
\tag{15}
\sum_{i=1}^{n} (a^{(1)} | x^{(i)}) (a^{(2)} | x^{(i)}) \cdots (a^{(n)} | x^{(i)}) &= (a^{(1)} \cdots a^{(n)})(x^{(1)} \cdots x^{(n)}).
\end{align}

Zwei weitere Identitäten gehen aus (13.) und (15.) durch die Substitutionen $(a^{(k)} | x) = (a^{(k)} q_{n-1})$, bzw.\ $(u | x^{(i)}) = (p_{\varrho} x^{(i)})$ hervor, sind also nach unserer Auffassung nicht als wesentlich verschieden zu betrachten. (14.) stellt den Produktsatz für Determinanten dar; wir wollen zeigen, wie man (13.) und (14.) (und dualistisch (15.) und (14.)) ersetzen kann durch das System der in geeigneter Weise gebildeten Produktsätze für Matrizen.

Der Produktsatz, einmal für $\varrho$-reihige, einmal für $(\varrho-1)$-reihige Matrizen gebildet, lautet:
\begin{align*}
(a^{(1)} \cdots a^{(\varrho)} | x p_{\varrho-1}) &= \sum_{i} (a^{(1)} | x^{(i)}) (a^{(2)} \cdots a^{(\varrho)} | p_{\varrho-1} x^{(i)}) \\
&= \sum (a^{(1)} | y) \cdots (a^{(\varrho)} | y) = (a^{(1)} \cdots a^{(\varrho)} | y \cdots y) = (a^{(1)} \cdots a^{(\varrho)} | p_{\varrho}),
\end{align*}
und daraus folgt das System von Produktsätzen, wo $\varrho = 2, 3, \ldots n$:
\begin{equation}\tag{16}
\sum_{i=1}^{\binom{n}{\varrho-1}} (-1)^{i} (a^{(1)} \cdots a^{(\varrho-1)} | x^{(i)} p_{\varrho-1}) (a^{(\varrho)} | x^{(i)}) = (a^{(1)} \cdots a^{(\varrho)} | p_{\varrho}).
\end{equation}

Für $\varrho = n$ geht (16.) in (13.) über; spezialisieren wir weiter: $p_{\varrho-1} = y p_{\varrho-2}$, $p_{\varrho-2} = y' p_{\varrho-3}$ usf., und wenden die Identitäten (16.) auf die Ausdrücke
\begin{equation*}
(a^{(1)} \cdots a^{(\varrho)} | x p_{\varrho-1}), \quad (a^{(1)} \cdots a^{(\varrho-1)} | x p_{\varrho-2} a^{(\varrho)})
\end{equation*}
usf.\ an, so gelangen wir von (13.) zu (14.). Die Identität (14.) ist also nach unserer Definition aus den Identitäten (16.) zusammengesetzt; oder das System (16.) ist den Identitäten (13.) und (14.) äquivalent. Das System (16.) erfüllt eine der für die allgemeinen Identitäten geforderten Bedingungen; es enthält die Reihe $p_{\varrho-1}$ als unzerlegbare Größe und zugleich in expliziter Form. Es ist das einzige System, das diese und nur diese Forderung erfüllt. Denn wir können aus den Reihen $a$,$x$,$p$ keine weiteren Determinanten oder Matrizenprodukte bilden; zwischen diesen können aber keine von (16.) unabhängigen Relationen bestehen, da sonst durch Elimination von $(a^{(1)} \cdots a^{(\varrho)} | p_{\varrho})$ eine Relation zwischen $\varrho$ ($\varrho < n$) Ausdrücken $(a^{(k)} | x)(a^{(1)} \cdots a^{(\varrho+1)} | q_{\varrho-1})$ entstände, nach (13.) aber mindestens $n+1$ solcher Ausdrücke in eine Relation eingehen müssen. Aus analogen Betrachtungen ergibt sich das (15.) und (14.) äquivalente System:
\begin{equation}\tag{17}
\sum_{i} (-1)^{i} (u^{(1)} \cdots u^{(n-\varrho-1)} | q_{n-\varrho} x^{(i)}) (u^{(n-\varrho)} | x^{(i)}) = (u^{(1)} \cdots u^{(n-\varrho)} | q_{n-\varrho+1}).
\end{equation}

Um von (16.) zu Identitäten zwischen beliebigen Reihen $A^{\varrho_{1}}$,$B^{\varrho_{2}}$,$\ldots$,$K^{\varrho_{k}}$,$p$ zu gelangen, bedenken wir, daß diese Identitäten mit (16.) identisch werden müssen, sobald wir auflösen: $A^{\varrho_{1}} = a^{(1)} \cdots a^{(\varrho_{1})}$, $B^{\varrho_{2}} = b^{(1)} \cdots b^{(\varrho_{2})}$ usf., und daß sie sich auch nur durch die durch (16.) bestimmten Operationen zu Schlußausdrücken zusammenfassen lassen, sobald die einzelnen Ausdrücke vom Typus der in (16.) eingehenden sind. Wir erhalten so, wenn wir noch
\begin{equation}\tag{18}
\varrho = \varrho_{1} + \varrho_{2} + \cdots + \varrho_{k} < n
\end{equation}
setzen:
\begin{equation}\tag{19}
\sum_{i=1}^{\binom{n}{\varrho_{1}}} (-1)^{i} (A^{\varrho_{1}} | x^{(i)})(B^{\varrho_{2}} \cdots K^{\varrho_{k}} | p_{\varrho-\varrho_{1}} x^{(i)}) = (A^{\varrho_{1}} B^{\varrho_{2}} \cdots K^{\varrho_{k}} | p_{\varrho}).
\end{equation}

Jede einzelne Summe (aber nicht schon ein Teil der Summe) enthält nun in der Tat die Reihen $A^{\varrho_{1}}$,$B^{\varrho_{2}}$,$\ldots$ als unzerlegbare Größe; denn sie erfüllt die dazu notwendigen und hinreichenden Bedingungen\footnote{Capelli, a.\ a.\ O.\ \S\ 23, gibt hinreichende Bedingungen:
\[
\sum_{\alpha=1}^{n} \frac{\partial F}{\partial x^{(k)}_{\alpha}} x^{(i)}_{\alpha} = 0, \quad \sum_{\alpha=1}^{n} \frac{\partial F}{\partial s^{(k)}_{\alpha}} \sigma^{(i)}_{\alpha} = 0 \quad (i \neq k; \, i,k = 1,2,\ldots \varrho-1)
\]
Daraus folgen am einfachsten die notwendigen und hinreichenden Bedingungen durch Anwendung des Poissonschen Klammerprozesses nach Wellstein, \textit{Von den Differentialgleichungen der projektiven Invarianten,} Math.\ Ann.\ 67 (1909) \S\ 3.}:
\begin{equation*}
\sum_{\alpha=1}^{n} \frac{\partial}{\partial x^{(k)}_{\alpha}} (\ldots) \cdot x^{(i)}_{\alpha} = 0, \qquad \sum_{\alpha=1}^{n} \frac{\partial}{\partial s^{(k)}_{\alpha}} (\ldots) \cdot \sigma^{(i)}_{\alpha} = 0 \quad (i \neq k).
\end{equation*}

Wir zeigen, daß sie alle Reihen auch explizit enthält. Denn setzen wir $B^{\varrho_{2}} \cdots K^{\varrho_{k}} = K^{\varrho-\varrho_{1}}$, so kommt nach (16.) und (10.):
\begin{align*}
(A^{\varrho_{1}} | x^{(i)})(K^{\varrho-\varrho_{1}} | p_{\varrho-\varrho_{1}} x^{(i)}) &= \text{nach (16.) und (10.)} \\
&= (-1)^{\varrho_{1}(\varrho-\varrho_{1})} (K^{\varrho-\varrho_{1}} A^{\varrho_{1}} | p_{\varrho}).
\end{align*}

Berechnen wir analog die übrigen Summen und setzen nachträglich, da die Reihen durch die Gewichtsindizes schon als verschieden charakterisiert sind, $B^{\varrho_{2}} = A^{\varrho_{2}}$, $C^{\varrho_{3}} = A^{\varrho_{3}}$, \ldots, so erhalten wir die gewünschten Identitäten:
\begin{equation}\tag{20}
(A^{\varrho_{1}} \cdots A^{\varrho_{k}} | p_{\varrho}) = \sum (-1)^{\varepsilon} (A^{\varrho_{1}} | x)(A^{\varrho_{2}} \cdots A^{\varrho_{k}} | p_{\varrho-\varrho_{1}} x) + (-1)^{\varepsilon'} (A^{\varrho_{1}} | p_{\varrho_{1}})(A^{\varrho_{2}} \cdots | p_{\varrho-\varrho_{1}}),
\end{equation}
wo $\varepsilon = (\varrho_{1} + \varrho_{2} + \cdots)(\varrho+1)$, und aus (17.) die dualistischen:
\begin{equation}\tag{21}
(U_{n-\varrho_{1}} \cdots | q_{n-\varrho}) = \sum (-1)^{\varepsilon} (U_{n-\varrho_{1}} | u)(\cdots | q_{n-\varrho+\varrho_{1}} u) + (-1)^{\varepsilon'} (U_{n-\varrho_{1}} | q_{n-\varrho_{1}})(\cdots | q_{n-\varrho+\varrho_{1}}),
\end{equation}
$\varepsilon = (\varrho_{1} + \varrho_{2} + \cdots + \varrho_{i-1}) \varrho_{i} + (\varrho_{i}+1)(n+\varrho)$.

\medskip
\textbf{Satz II.} \textit{Mit (20.) und (21.) ist die Gesamtheit der unzerlegbaren Identitäten zwischen den Reihen $A$,$x$,$p$, bzw.\ $A$,$u$,$q$ erschöpft, und zugleich die Gesamtheit derjenigen unzerlegbaren Identitäten, die explizite Darstellung ohne Vornahme der Substitution (11.) zulassen.}
\medskip

Der erste Teil des Satzes ergibt sich aus den Betrachtungen zur Ableitung von (20.) und (21.); der zweite Teil aus der Unmöglichkeit einer Identität zwischen zwei Symbolreihen:
\begin{equation*}
(q_{\varrho} A^{\varrho} | B^{\tau} p_{\tau}) = \varepsilon_{1} (A^{\varrho} B^{\tau} | q_{n-\varrho} p_{\tau}) + \varepsilon_{2} (q_{\varrho} B^{\tau} | A^{n-\varrho} p_{\varrho}),
\end{equation*}
wo $\tau < \varrho$ und keine der Reihen vom Gewicht $1 \cdot (n-1)$ ist. (Andere Annahmen über $\varrho$,$\tau$,$\sigma$ lassen sich durch Vertauschen der Reihenbezeichnungen auf diese zurückführen.) Diese Unmöglichkeit folgt z.\ B.\ durch Entwicklung (8.) bei der Spezialisierung $q_{\varrho} = I \cdot M$,$q_{\tau} = I \cdot N$, wenn man $I_{1} = 1$;$M^{\varrho+1} = 1$;$A^{\varrho+\tau-1} = 1$, alle anderen Elemente der Reihen $I$,$M$,$A$ gleich Null setzt, während $N$ und $B$ beliebig sind. Da nun Identitäten zwischen beliebigen Reihen durch Auflösen einer Reihe $p_{\varrho}$ in $y^{(1)} y^{(2)} \cdots y^{(\varrho)}$ in solche übergehen, die aus (20.) zusammengesetzt sind, ist der zweite Teil der Behauptung fertig bewiesen, sobald sich (20.) aus Identitäten mit nur zwei Symbolreihen erzeugen läßt.

Man erhält in der Tat aus
\begin{equation*}\tag{20$^{*}$}
(A^{\varrho} B^{\tau} | p_{\varrho+\tau}) = \varepsilon (A^{\varrho} q_{n-\varrho} | B^{\tau} p_{\tau}) + (A^{\varrho} | p_{\varrho}),
\end{equation*}
wo $p_{\varrho+\tau} = A^{\varrho} B^{\tau}$, durch die Spezialisierung $A^{\varrho} = B' B''$, d.h.\ durch
\begin{equation*}
(B' B'' q_{\varrho} | \ldots) = \varepsilon (B' q_{\varrho} q_{\varrho+1} | B'' \ldots) + (B' | \ldots),
\end{equation*}
eine Identität (20.) in drei Symbolreihen; und durch weitere Spezialisierungen von $B'' = C' C''$ usf.\ die allgemeinen Identitäten (20.).

\section*{\S\ 4. Die Zerlegungsidentitäten.}
\addcontentsline{toc}{section}{\S\ 4. Die Zerlegungsidentitäten}

Wir betrachten nun den allgemeinsten Fall der Identitäten zwischen beliebigen Symbol- und Variabelnreihen, in deren Schlußausdrücken ohne Vornahme der Substitution (11.) eine Reihe $p_{\varrho}$ in ihre einzelnen Reihen $y$ zerlegt auftritt, und bezeichnen diese Identitäten als ``Zerlegungsidentitäten.'' Im Anschluß an die Bemerkung am Schluß des vorigen Paragraphen bestimmen sie vorerst für den Fall nur zweier Symbolreihen, entwickeln also den Ausdruck $(q_{\varrho} A^{\varrho} | B^{\tau} p_{\tau})$. Wir setzen dabei:
\begin{equation}\tag{22}
p_{\varrho} = [y^{(1)} \cdots y^{(\sigma)} \cdot y^{(\sigma+1)} \cdots y^{(\varrho)}]_{\varrho}, \quad p_{\tau} = [y^{(\sigma+1)} \cdots y^{(\varrho)} \cdot y^{(\varrho+1)} \cdots y^{(\varrho+\tau)}]_{\tau}
\end{equation}
und
\begin{equation*}
\varrho - \sigma = \tau, \quad \sigma + \tau = \varrho.
\end{equation*}

Inbezug auf die Zahlen $\varrho$,$\sigma$,$\tau$ sind vier Fälle zu unterscheiden:
\begin{center}
\begin{tabular}{ll@{\qquad}ll}
1) $\varrho + \sigma < n$ & $\tau > \sigma$ & 2) $\varrho + \sigma < n$ & $\tau > \sigma$ \\
3) $\varrho + \sigma \geq n$ & $\tau > \sigma$ & 4) $\varrho + \sigma < n$ & $\tau \leq \sigma$,
\end{tabular}
\end{center}
von denen wir den ersten herausgreifen wollen, um danach die übrigen kurz zu charakterisieren. Nach (20.), (10.) und (9.) kommt, wenn wir die Reihen $y$ geeignet wählen:
\begin{equation}\tag{23}
(q_{\varrho} A^{\varrho} | y^{\sigma} B^{\tau} y^{\varrho-\sigma} p_{\tau}) = (q_{\varrho} A^{\varrho} y^{\sigma} | B^{\tau} y^{\varrho-\sigma} p_{\tau}) + (q_{\varrho} A^{\varrho} | y^{\sigma} B^{\tau} y^{\varrho-\sigma} p_{\tau}).
\end{equation}

Das letzte Glied in (23.) hat nun genau die Form des ursprünglichen Ausdrucks --- es ist nur $\sigma$ durch $\sigma-1$, $\tau$ durch $\tau-1$ ersetzt --- läßt also wieder die Gleichung (23.) zu. Wir können also, da $\tau < \varrho$, die Operation (23.) $\tau$-mal anwenden und gelangen unter Berücksichtigung von (10.) zu der Relation:
\begin{equation}\tag{24}
(q_{\varrho} A^{\varrho} | y^{\sigma} B^{\tau} y^{\varrho-\sigma} p_{\tau}) = \sum_{i=0}^{\tau-1} (-1)^{\varepsilon_{i}} (q_{\varrho} A^{\varrho} y^{\sigma+i} | B^{\tau} y^{\varrho-\sigma-i} p_{\tau}) + (-1)^{\varepsilon_{\tau}} (q_{\varrho} A^{\varrho} y^{\sigma+\tau} | B^{\tau} p_{\varrho-\sigma-\tau}).
\end{equation}

Ein jedes Glied unter dem Summenzeichen ist nun wieder vom Typus des ursprünglichen Ausdrucks; es ist $\varrho$ durch $\varrho-1$, $\sigma$ durch $\sigma+\tau_{i}+1$, $\tau$ durch $\tau-1$ ersetzt, die Bedingung 1) also noch stärker erfüllt. Wir können deshalb, da $\tau < \varrho$, die Operation (24.) $\tau$-mal anwenden und erhalten die gewünschte Zerlegungsidentität:
\begin{multline}\tag{25}
(q_{\varrho} A^{\varrho} | B^{\tau} p_{\tau}) = \sum_{i_{1}=0}^{\tau-1} \sum_{i_{2}=i_{1}+1}^{\tau-1} \cdots \sum_{i_{\sigma}=i_{\sigma-1}+1}^{\varrho-\sigma} \varepsilon_{i_{1} \cdots i_{\sigma}} \\ \cdot (q_{\varrho} A^{\varrho} y^{\sigma} \cdots y^{\sigma+i_{\sigma}} | B^{\tau} y^{\varrho-\sigma-i_{\sigma}} \cdots y^{\varrho-\sigma} p_{\tau}).
\end{multline}

Für den Modul $\varepsilon$ ergibt sich aus (24.):
\begin{equation}\tag{26}
\varepsilon_{i} = (-1)^{M_{i}}, \quad M = \sum (\varrho - \beta + 1)(i_{\beta} + \beta) + (\varrho - \sigma)(\sigma + \tau_{i} + \varrho) + (\sigma + \tau + \varrho)(n+1),
\end{equation}
wo $\tau_{0} = 0$ zu setzen ist. Das Summenzeichen in (25.) ist so zu verstehen, daß jeder Kombination $i_{1}, i_{2}, \ldots, i_{\sigma-1}$, wo $1 < i_{1} < \cdots < i_{\sigma-1}$, alle Werte $i_{\sigma}$, für die $i_{\sigma-1} < i_{\sigma} < \varrho - \sigma$, zugefügt werden. Die einzelnen Summen in (25.) genügen, wie man sich unter Berücksichtigung von (26.) leicht überzeugt, den (19.) analogen Differentialgleichungen:
\begin{equation*}
\sum_{\alpha} \frac{\partial}{\partial x^{(k)}_{\alpha}} (\ldots) \cdot x^{(i)}_{\alpha} = 0,
\end{equation*}
hängen also nur von der Reihe $p_{\varrho}$ ab. Die Zerlegungsidentität ist also eine unzerlegbare Identität zwischen den Reihen $A$,$B$,$q$,$p$; auf die explizite Darstellung kommen wir im nächsten Paragraphen zurück.

Im Fall 2) läßt sich die Operation (23.) ebenfalls $\tau$-mal anwenden, die Operation (24.) aber bricht nach dem $\varrho$-ten Male ab. Es ist also im Schlußausdruck von (25.) das erste Summenzeichen zu ersetzen durch $\sum_{i_{1}=0}^{\varrho-\sigma-1}$.

Im Fall 3) und 4) läßt sich die Operation (23.) nur $\sigma$-mal ausführen; an Stelle von (24.) tritt dann die Relation:
\begin{equation}\tag{28}
(q_{\varrho} A^{\varrho} | y^{\sigma} B^{\tau} y^{\varrho-\sigma} p_{\tau}) = \sum_{i=0}^{\sigma-1} (-1)^{\varepsilon_{i}} (q_{\varrho} A^{\varrho} y^{\sigma-i} | B^{\tau} y^{\varrho-\sigma+i} p_{\tau}),
\end{equation}
die $(\varrho-\sigma)$-malige Wiederholung von (28.), wobei das Summenzeichen $\sum_{i_{\sigma}=i_{\sigma-1}+1}^{\tau_{\sigma}-1}$ auftritt, --- wie die $(\varrho-\tau_{i}+1)$-malige Anwendung von (23.) auf die Glieder unter dem Summenzeichen in (28.) zeigt ---, gibt alle Glieder der dem Index $\varepsilon = \tau - \sigma$ entsprechenden Einzelsumme von (25.) mit entsprechendem Vorzeichen. Dann aber geht Fall 3) und 4) in Fall 1) und 2) über, denn die $\sigma$,$\tau$ entsprechenden Zahlen werden: $\sigma' = \sigma - \tau_{i-1} + \tau - \varrho$, $\tau' = \tau - i_{\sigma-1} = 0$; und bei Fortsetzung des Verfahrens wird $\tau' < \sigma'$. Es ist also im Schlußausdruck von (25.) das erste Summenzeichen zu ersetzen durch $\sum_{i_{1}=0}^{\tau-\sigma}$.

Analog ergibt sich aus (21.) die dualistische Zerlegungsidentität, bei der eine Reihe $q_{\varrho}$ in die einzelnen Reihen $v^{(1)}, \ldots, v^{(\varrho)}$ zerlegt wird. Setzen wir
\begin{equation}\tag{29}
p_{\varrho} = [v^{(1)} \cdots v^{(\sigma)} \cdot v^{(\sigma+1)} \cdots v^{(\varrho)}]_{\varrho}, \quad q_{\varrho} = [v^{(\sigma+1)} \cdots v^{(\varrho)} \cdot v^{(\varrho+1)} \cdots v^{(\varrho+\tau)}]_{n-\varrho},
\end{equation}
so kommt:
\begin{equation}\tag{30}
(q_{\varrho} A^{\varrho} | B^{\tau} p_{\tau}) = \sum_{\sigma=0}^{\varrho} \sum_{\tau=0}^{\varrho-\sigma} \varepsilon_{\sigma,\tau} (q_{\varrho} A^{\varrho} v^{\sigma} | B^{\tau} v^{\varrho-\sigma} p_{\tau}),
\end{equation}
wo der Summationsindex $\sigma$ vom größeren der unteren Werte bis zum kleineren der oberen läuft. Alle Glieder einer Einzelsumme in (25.) und (30.) sind Polaren einer einzelnen Form, entstanden durch die Polarprozesse $\Pi = (\frac{\partial}{\partial t_{1}} \cdots \frac{\partial}{\partial t_{\sigma}})(i_{1}, \ldots, i_{\sigma}) (\frac{\partial}{\partial y_{1}} \cdots \frac{\partial}{\partial y_{\varrho}})$. Um dies zum Ausdruck zu bringen, bezeichnen wir mit $\Pi$ ein Aggregat solcher Polaroperationen, multipliziert mit Konstanten, und erhalten dann an Stelle von (25.) und (30.) die Relation:
\begin{equation}\tag{31}
\Pi(q_{\varrho} A^{\varrho} | B^{\tau} p_{\tau}) = \sum \varepsilon \cdot \Pi(q_{\varrho} A^{\varrho} y^{\sigma} | B^{\tau} y^{\varrho-\sigma} p_{\tau}).
\end{equation}

\section*{\S\ 5. Die Zerlegungsidentitäten in expliziter Form.}
\addcontentsline{toc}{section}{\S\ 5. Die Zerlegungsidentitäten in expliziter Form}

Spezialisieren wir in (25.) im Falle 4) $\tau$ zu $\sigma + \varrho$, so kommt:
\begin{equation}\tag{32}
(A^{\varrho} y^{1} \cdots y^{\varrho} | B^{\tau} p_{\tau}) = \sum_{\alpha=1}^{\varrho} \varepsilon_{\alpha} (q_{\varrho} y^{1} \cdots y^{\varrho} | A^{\varrho} y^{\alpha+1} \cdots y^{\varrho+\alpha} p_{\varrho}),
\end{equation}
wo $\varepsilon = (-1)^{M}$, $M = \varrho_{\alpha_{1}} + \varrho_{\alpha_{2}} + \cdots$ wird. Dies ist eine allgemeinere Form des Produktsatzes für Matrizen als (16.) und (17.) und ergibt sich durch Laplacesche Entwicklung der Determinante der Produktenmatrix:
\begin{equation*}
\sum (v^{(1)} | y^{(1)}) \cdots (v^{(\varrho)} | y^{(\varrho)}) (a^{(1)} | y^{(\varrho+1)}) \cdots (a^{(\tau)} | y^{(\varrho+\tau)}).
\end{equation*}
Der allgemeinere Produktsatz ist also aus den Identitäten (20.), bzw.\ (21.) zusammengesetzt, und das Herausgreifen der speziellen Form (16.), (17.) ist dadurch erklärt.

Die Relation (32.) gibt nun das Mittel zur expliziten Darstellung der Zerlegungsidentitäten. Wir betrachten dazu die in (31.) auftretenden Formen als Grundformen; d.\ h.\ wir führen die Substitution (11.) aus, die an der expliziten Darstellung in den Reihen $A$,$B$ nichts ändert, da nach (8.) die neu eingeführten Reihen $R$ sich durch die Reihen $A$,$B$ selbst und nicht durch deren Teilreihen $a^{(1)} a^{(2)} \ldots$,$b^{(1)} b^{(2)} \ldots$ ausdrücken lassen. Sei also
\begin{equation*}
(G_{n} | R^{\varrho_{1}} B^{\tau_{1}}) = (R^{\varrho_{1}} | B^{\tau_{1}}) \cdots,
\end{equation*}
so kommt für die dem Index $\varepsilon$ entsprechende Einzelsumme von (25.):
\begin{multline}\tag{34}
\sum \varepsilon_{\alpha} (q_{\varrho-\alpha} p_{\varrho} y^{\alpha+1} \cdots y^{\varrho} | B^{\tau-\alpha} y^{\varrho-\alpha+1} \cdots y^{\varrho}) \\ = \sum \varepsilon_{\alpha} (R^{\varrho_{1}} \cdots | q_{\varrho} y^{\alpha+1} \cdots y^{\varrho}) (R^{\varrho_{i}} | y^{\varrho-\alpha+1} \cdots y^{\varrho} p_{\varrho}) \\ = \varepsilon (R^{\varrho_{i}} | q_{\varrho} p_{\varrho}) \cdots.
\end{multline}

Damit ist nun in der Tat ein expliziter Schlußausdruck gegeben, und die symmetrische Form, in der $q_{\varrho}$ und $p_{\varrho}$ eingehen, zeigt, daß (30.) zum gleichen Schlußausdruck führt, der also unabhängig ist von der ursprünglichen Zerlegung. Ein Unterschied zwischen (25.) und (30.) besteht nur, solange die Reihen $y$ und $v$ selbständige Bedeutung haben, die Zerlegungsidentität also nach unserer Definition in eine zerlegbare Identität übergeht.

Um zu Identitäten mit mehr Symbolreihen zu gelangen, hat man nach \S\ 3 (Schluß) nur eine Reihe $q_{\varrho}$ in $q_{\varrho'} C^{\varrho''}$, bzw.\ $p_{\varrho}$ in $p_{\varrho'} C^{n-\varrho''}$, aufzulösen. Die dann entstehenden Ausdrücke:
\begin{equation*}
(q_{\varrho'} C^{\varrho''} A^{\varrho} | B^{\tau} p_{\tau}) \quad \text{bzw.} \quad (q_{\varrho} A^{\varrho} | B^{\tau} p_{\varrho'} C^{n-\varrho''})
\end{equation*}
sind nun genau vom Typus der bei zwei Symbolreihen auftretenden, lassen also bei einer (33.) entsprechenden Substitution wieder explizite Darstellung zu, so daß sich durch Wiederholung explizite Darstellung bei beliebig vielen Reihen ergibt. Bemerkt mag noch werden, daß je die erste und letzte Summe in (25.) bzw.\ (30.) die explizite Darstellung ohne Anwendung von (33.) allein durch die Formel (32.) ergibt, oder sich überhaupt auf ein einziges, alle Reihen explizit enthaltendes Glied beschränkt. Die durch Auflösung von $q_{\varrho}$ in $q_{\varrho'} C^{\varrho''}$ usf.\ aus (32.) entstehende Relation kann als der Produktsatz in seiner allgemeinsten Form aufgefaßt werden.

Zusammenfassend werden wir zeigen:

\medskip
\textbf{Satz III:} \textit{Mit den Identitäten}
\begin{equation}\tag{36}
(q_{\varrho} A^{\varrho} | B^{\tau} p_{\tau}) = \sum \varepsilon (R^{\varrho_{i}} | q_{\varrho} p_{\varrho}) (R^{\varrho_{i}} | \ldots),
\end{equation}
\textit{wo die Reihen $R$ durch (33.) bestimmt sind, und mit den durch Auflösung von $q_{\varrho}$,$p_{\varrho}$ in $q_{\varrho'} C^{\varrho''}$ usf.\ daraus entstehenden ist die Gesamtheit der unzerlegbaren Identitäten erschöpft.}
\medskip

In der Tat sind die entstehenden Identitäten die allgemeinsten ihrer Art, da die einzelnen Reihen in bezug auf Gewicht und Anzahl keiner Beschränkung mehr unterliegen, wie dies noch bei (20.), (21.) der Fall war. Es existieren aber auch keine weiteren Identitäten zwischen diesen Reihen; denn bei Auflösung einer Reihe $p_{\varrho}$ in $y^{(1)} \cdots y^{(\varrho)}$ sind die allgemeinsten Identitäten aus (20.), also nach \S\ 3 (Schluß) aus (20$^{*}$.) zusammengesetzt; und zwar ist die im Anschluß an (16.) diskutierte Zusammensetzung genau die in \S\ 4 durchgeführte. Den Übergang zu beliebig vielen Reihen liefert aber, wie die Diskussion von (20$^{*}$.) zeigt, die Anwendung immer derselben Operation auf die durch Auflösung von $q_{\varrho}$ in $q_{\varrho'} C^{\varrho''}$ usf.\ entstehenden Ausdrücke. Es folgt daraus auch, daß der direkte Übergang von (20.) an Stelle von (20$^{*}$.) nichts Neues liefert; dies läßt sich rechnerisch leicht bestätigen, indem man einmal eine Reihe $q_{\varrho}$ in (25.) spezialisiert, das andere Mal nach der Methode von \S\ 4 den Übergang von (20.) macht; es entstehen beide Male die gleichen Summen, die durch Anwendung des allgemeinen Produktsatzes (siehe oben) in die aus (36.) sich ergebenden Identitäten übergehen. Die Identitäten (20.), (21.) entsprechen den Spezialfällen $\sigma = 1$,$\varrho = 1$; es treten dann nur je zwei Einzelsummen auf, also einer obigen Bemerkung entsprechend explizite Darstellung ohne die Substitution (33.), was mit dem früheren übereinstimmt.

Zum Schluß seien noch zwei Spezialfälle der Zerlegungsidentität kurz erwähnt. Setzen wir in (31.): $\tau = \sigma$, $\varrho = n-\sigma$ und bezeichnen die Reihe $p_{\varrho}$ als $p_{\varrho} = f$, so kommt:
\begin{equation}\tag{37}
(A^{\varrho} f | B^{\tau} p_{\tau}) - (A^{\varrho} | f)(B^{\tau} | p_{\tau}) = \sum \Pi (q_{\varrho} A^{\varrho} B^{\tau} | f y^{\sigma} p_{\varrho}).
\end{equation}
Das ist die Clebschsche Formel\footnote{Clebsch, a.a.O., S.~25--32. Clebsch weist nach, daß die einzelnen Glieder der rechten Seiten nur von den Reihen $A$,$B$ abhängen; die explizite Darstellung würde sich durch Anwendung von (32.) auf seine Ausdrücke ergeben.} zur Entwicklung einer Form mit zwei kogredienten Reihen $p_{\varrho}$,$p'_{\varrho}$ nach Elementarkovarianten. Die zu einer Form $(A^{\varrho} | p_{\varrho})^{m} (B^{\tau} | p_{\tau})^{m'}$ gehörenden Elementarkovarianten sind dann die folgenden:
\begin{equation}\tag{38}
\Pi_{\sigma} = (q_{\varrho} A^{\varrho} y^{\sigma} | p_{\varrho} y^{\varrho-\sigma}) (B^{\tau} | p_{\tau}) \qquad (\sigma = 0, 1, \ldots m, n).
\end{equation}

Setzen wir zweitens: $\tau = \varrho - \lambda$, $\sigma = n-\varrho-1$; $B^{\tau} = A^{\varrho}$, so kommt:
\begin{equation}\tag{39}
(A^{\varrho} | q_{\varrho} p_{\varrho}) = \sum \Pi (A^{\varrho} y^{\sigma} | q_{\varrho} y^{\varrho-\sigma} p_{\varrho}).
\end{equation}
Es sind also, wie in \S\ 2 bemerkt, die die Symbolreihen charakterisierenden Bedingungen (12.) für ungeraden Defekt $\lambda$ eine Folge der Bedingungen von höherem Defekt. Für eine Linearform $(A^{1} | p_{1}) = (B^{1} | p_{1})$ ergibt sich aus (39.), daß sich ihre irreduzibeln invarianten Bildungen, die in den Koeffizienten quadratisch sind, auf solche von geradem Defekt reduzieren; das stimmt mit der bekannten Tatsache überein, daß eine Linearform im binären und ternären Gebiet keine invarianten Bildungen besitzt.

Bemerkt sei, daß die allgemeinen Identitäten ursprünglich abgeleitet sind unter der Annahme, daß die einzelnen Reihen den Bedingungen (12.) genügen. Da indes diese einzelnen Reihen nur linear in die Identitäten eingehen, so gelten letztere auch für die allgemeineren, den Bedingungen (12.) nicht genügenden Reihen.

\section*{\S\ 6. Explizite Darstellung der invarianten Bildungen. Faltungs- und Differentiationsprozesse.}
\addcontentsline{toc}{section}{\S\ 6. Explizite Darstellung der invarianten Bildungen}

Die Resultate der vorigen Paragraphen erlauben uns, den in \S\ 2 ausgesprochenen Satz von der symbolischen Darstellbarkeit (1.\ Fundamentalsatz) durch Zufügung der expliziten Darstellbarkeit zum Abschluß zu bringen, indem wir zeigen:

\medskip
\textbf{Satz IV:} \textit{Alle invarianten Bildungen beliebiger Grundformen lassen sich explizit durch Matrizenprodukte darstellen; d.\ h.\ in die Schlußausdrücke gehen außer Variablenreihen $p_{\varrho}$ nur die Reihen $S^{\varrho}$,$\ldots$,$T^{\tau}$ der ursprünglichen Formen oder die durch Substitution (9.) oder (11.) aus ihnen abgeleiteten Reihen $P$,$Q$,$R$ ein.}
\medskip

Zum Beweise nehmen wir an, eine invariante Bildung möge außer von Variablenreihen von den Reihen $S^{\varrho}$,$\ldots$,$T^{\tau}$ abhängen; und zwar seien diese Reihen genügend weit in einzelne Reihen: $S^{\varrho} = s^{(1)} \cdots s^{(\varrho)}$, \ldots, $T^{\tau} = t^{(1)} \cdots t^{(\tau)}$ aufgelöst, um eine Darstellung durch Determinanten zu ermöglichen. Die Reihen seien in einer an sich willkürlichen, aber bestimmt festgelegten Reihenfolge $S^{\varrho}$,$\ldots$,$T^{\tau}$ geordnet. Wir greifen nun ein solches Aggregat von Determinanten heraus, das von der ersten Gesamtreihe $S^{\varrho} = s^{(1)} \cdots s^{(\varrho)}$, nicht nur von einzelnen der Teilreihen, abhängt. Diese Abhängigkeit ist aber nach \S\ 3--\S\ 5 eine Folge der allgemeinen Identitäten. Lösen wir nun eine Variablenreihe $p_{\varrho}$ in die Reihen $y$,$v$, oder eine der späteren Reihen $T$ in die einzelnen Reihen $t$, bzw.\ $\tau$ auf, so sind die allgemeinsten Identitäten aus (20.), (21.) zusammengesetzt. Diese letzteren Identitäten führen aber immer zu einem die Reihe $S^{\varrho} = s^{(1)} \cdots s^{(\varrho)}$ explizit enthaltenden Ausdruck, der linken Seite von (20.), (21.); man überzeugt sich nämlich durch (19.), daß erst die vollständige rechts stehende Summe --- nicht aber ein Teil derselben, der einem Gliede von (20$^{*}$.) entspricht --- die Eigenschaft besitzt, nur von $S^{\varrho}$ abzuhängen. Bei Fortsetzung des Verfahrens bleiben alle einmal explizit eingehenden Reihen auch explizit erhalten; die Anwendung der Identitäten kann nur möglicherweise ein Zusammenfassen mehrerer Reihen zu einer einzigen, d.\ h.\ die Substitution (9.), verlangen. Ist schließlich nur noch eine einzige Reihe $T^{\tau}$ auf die explizite Form zu bringen, so sind zwei Fälle möglich. Es kann noch eine freie, d.\ h.\ nicht durch Substitution (9.) mit andern Reihen verbundene, Variabelnreihe auftreten; durch deren Zerlegung in Teilreihen $y$ oder $v$ läßt sich das Verfahren zum Abschluß bringen; es entstehen, wie in (31.), Polaren einer oder mehrerer, alle Reihen explizit enthaltenden Formen. Tritt keine freie Variabelnreihe mehr auf, so erkennen wir nach der Art, wie sie entstanden sind, in der Gesamtheit der Ausdrücke, die die einzelnen Reihen $t$ oder $\tau$ enthalten, aber von der Reihe $T^{\tau}$ abhängen, die Glieder einer Zerlegungsidentität; diese aber lassen nach \S\ 5 unter Anwendung der Substitution (11.) stets die explizite Darstellung in allen Reihen zu. Damit ist der obige Satz allgemein bewiesen. Bemerkt mag noch werden, daß dann, wenn nur zwei Symbolreihen auftreten, die Substitution (11.) nie eintritt; es können nämlich --- wegen $\varrho$,$n-\varrho$,$\tau$,$n-\tau < n$ --- nur dann keine Variabelnreihen eingehen, wenn die beiden Symbolreihen in einer einzigen Determinante vereinigt, also explizit auftreten; sobald aber Variabelnreihen eingehen, zeigen (34.) und (33.), daß die Substitution (11.) überflüssig wird.

Aus dem eben Gesagten folgt, daß es zur Erzeugung aller invarianten Bildungen genügt, die Symbolreihen unter Zufügung von Variablenreihen zu allen möglichen Matrizenprodukten zu vereinigen. Das allgemeinste Matrizenprodukt ist nun von der Form:
\begin{equation}\tag{40}
(P^{\varrho_{1}} \cdots P^{\varrho_{k}} | Q_{\tau_{1}} \cdots Q_{\tau_{l}}), \qquad \varrho_{i} \gtrless \tau_{j},
\end{equation}
wo die $P$ und $Q$ Reihen der ursprünglichen Formen sein können, oder durch eine Reihenfolge von Substitutionen (9.) oder (11.) aus den ursprünglichen Reihen hervorgegangen, oder endlich Variabelnreihen. Die Ausdrücke (40.) aber lassen sich erzeugen durch sukzessive Vereinigung je zweier Symbolreihen zu einem Matrizenprodukt, unter Zuhilfenahme von Variablenreihen; denn man erhält aus
\begin{equation*}
(G_{n} | R^{\varrho} Q_{\tau}) = (R^{\varrho} A^{\varrho} | Q_{\tau} P_{\varrho})
\end{equation*}
durch Vereinigung mit der Reihe $Q_{\tau_{2}}$, einen Ausdruck
\begin{equation*}
(G_{n} | R^{\varrho} Q_{\tau_{1}} Q_{\tau_{2}}) = (R^{\varrho} A^{\varrho} | Q_{\tau_{1}} Q_{\tau_{2}} P_{\varrho}),
\end{equation*}
also durch Fortsetzung des Verfahrens einen Ausdruck (40.). Sind dabei $P$ und $Q$ nicht die Reihen der ursprünglichen Formen, so zeigt (9.) und (11.) in Verbindung mit dem eben Gezeigten, daß sie gleichfalls durch eine Reihenfolge von Vereinigungen je zweier Symbolreihen entstanden sind. Es folgt daraus:

\medskip
\textbf{Satz V:} \textit{Der Prozeß zur Erzeugung aller invarianten Bildungen, der Faltungsprozeß, besteht in der sukzessiven Vereinigung je zweier Symbolreihen zu einem Matrizenprodukt, unter Zuhilfenahme von Variablenreihen.}
\medskip

Aus zwei Symbolreihen $S^{\varrho}$,$T^{\tau}$, wo $\varrho > \tau$, lassen sich nun die folgenden Matrizenprodukte bilden:
\begin{align}\tag{41}
&(G_{n-\varrho-\lambda} S^{\varrho} T^{\tau} | p_{\varrho+\lambda} p_{\tau-\lambda}), \\ \tag{42}
&(G_{n-\varrho-\lambda-1} S^{\varrho} T^{\tau} | p_{\varrho+\lambda+1} p_{\tau-\lambda-1}), \qquad (\mu = 0, 1, \ldots \tau, \varrho-\varrho) \\ \tag{43}
&(G_{n-\varrho} S^{\varrho} T^{\tau} | q_{n-\varrho} p_{\tau}) \qquad (\lambda = 1, 2, \ldots \tau, n-\varrho).
\end{align}

Aus (31.) ergeben sich aber für (42.) und (43.) die Werte:
\begin{align*}\tag{42$^{*}$}
&(G_{n-\varrho-\lambda-1} S^{\varrho} T^{\tau} | p_{\varrho+\lambda+1} p_{\tau-\lambda-1}) = \varepsilon \sum_{i=1}^{\lambda} (G_{n-\varrho-i} S^{\varrho} | p_{\varrho+i} p_{\tau-i}), \\ \tag{43$^{*}$}
&(G_{n-\varrho} S^{\varrho} T^{\tau} | q_{n-\varrho} p_{\tau}) = \varepsilon \sum_{i=1}^{\tau} (G_{n-\varrho-i} S^{\varrho} | p_{\varrho+i} p_{\tau-i}) + \varepsilon (T^{\tau} | p_{\tau}),
\end{align*}
$(\lambda = 1, 2, \ldots \varrho-1)$. Es lassen sich also die Formen (42.) linear ausdrücken durch (41.) und Polaren von (41.), die Formen (43.) durch Polaren der Grundform und der Formen (41.). Ebenso ergibt sich aus (31.) die lineare Abhängigkeit der Formen (41.) von (42.) und Polaren von (42.). Danach sind aber nach der Clebschschen Definition\footnote{Clebsch, a.a.O.\ \S\ 7.} die Formen (42.) den Formen (41.) äquivalent, während (43.) auf die Grundform zurückführt. Der erzeugende Prozeß ist also gleichberechtigt durch (41.) oder (42.) gegeben; wir wollen den in bezug auf die Zahl $\lambda$ einfacheren Prozeß (41.) als Faltungsprozeß definieren. Die Zahl $\lambda$, der Defekt der Faltung (vgl.\ \S\ 2) gibt die Anzahl der zum Determinantenprodukt fehlenden Reihen an, und zwar in demjenigen Matrizenprodukt, das bei gegebener Faltung die Maximalzahl von Reihen enthält. Da $\lambda$ nur bis zum kleineren der beiden Werte $\tau$,$n-\varrho$ läuft, wo $\varrho > \tau$, so ist:
\begin{equation}\tag{44}
\lambda \leq \frac{n}{2}, \quad n \text{ gerade}; \qquad \lambda \leq \frac{n-1}{2}, \quad n \text{ ungerade}.
\end{equation}

Die Zurückführung von (42.) und (43.) auf (41.) gilt auch dann noch, wenn durch weitere Faltung die Reihen $p$,$q$ in Symbolreihen übergegangen sind; denn nach (36.) kommt:
\begin{equation*}
(A^{\varrho} B^{\tau} C^{\sigma} | p_{\varrho} p_{\tau} p_{\sigma}) = \sum \varepsilon (G_{n} | R^{\varrho_{i}} \ldots) \cdot (R^{\varrho_{i}} | \ldots),
\end{equation*}
wo die Reihen $R$ nach (33.) aus $S$ und $T$ hervorgegangen sind. Man erhält also diese nur Symbolreihen enthaltenden Formen durch Faltung von $(G_{n-\varrho} A^{\varrho} B^{\tau} | p_{\varrho} p_{\tau})$ bzw.\ $(G_{n-\tau} B^{\tau} A^{\varrho} | p_{\tau} p_{\varrho})$ --- je nachdem $\varrho-1 \gtrless \tau$ --- mit der Grundform und den Formen (41.).

Aus dem symbolischen Faltungsprozeß gewinnt man in bekannter Weise die invarianten Differentiationsprozesse, wenn man nur die Symbolreihen $S^{\varrho}$,$T^{\tau}_{n-\tau}$ durch die Reihen von Differentiationssymbolen $\frac{\partial}{\partial s}$,$\frac{\partial}{\partial t}$ ersetzt, wobei bekanntlich jedem Produkte $\frac{\partial f}{\partial s} \cdot \frac{\partial \varphi}{\partial \sigma}$ die reale Bedeutung $\sum_{\alpha} \frac{\partial f}{\partial s_{\alpha}} \frac{\partial \varphi}{\partial \sigma_{\alpha}}$ zukommt. Da die Reihen $\frac{\partial f}{\partial s}$ den Reihen $S^{\varrho}$ sind, genügen sie denselben Identitäten wie diese, also speziell auch den zwischen (42.) und (42$^{*}$.), (43.) und (43$^{*}$.) bestehenden.

Zusammenfassend erhalten wir:

\medskip
\textbf{Satz VI:} \textit{Alle invarianten Bildungen beliebiger Grundformen entstehen aus diesen durch wiederholte Anwendung des symbolischen Faltungsprozesses (41.) oder auch der invarianten Differentiationsprozesse:}
\begin{equation}\tag{45}
D_{\lambda} = (G_{n-\varrho-\lambda} \frac{\partial}{\partial s^{(1)}} \cdots \frac{\partial}{\partial s^{(\varrho)}} | p_{\varrho+\lambda} p_{\tau-\lambda}) \qquad (\varrho > \tau, \lambda = 1, 2, \ldots \tau, n-\varrho).
\end{equation}
\medskip

Es ist noch zu bemerken, daß bei jeder Faltung (41.) das Gesamtgewicht der Form, d.\ h.\ die Summe der Gewichte der einzelnen Variabelnreihen (vgl.\ \S\ 1), abnimmt um die positive Zahl $2(\lambda^{2} + \lambda(\varrho-\tau))$. Daraus folgt, unabhängig vom allgemeinen Endlichkeitsbeweis, die Endlichkeit der Formenreihe, d.\ h.\ der Gesamtheit der in den Koeffizienten linearen invarianten Bildungen einer gegebenen Grundform\footnote{Vgl.\ Clebsch, a.\ a.\ O.\ \S\ 11 u.\ 12.\ Endlichkeit des Systems von Elementarkovarianten.}.

Bemerkt sei weiter, daß Satz VI auch noch gilt für Formen, die den Bedingungen (12.) nicht genügen. Es läßt sich nämlich jeder Faktor von (3.): $(S^{\varrho} | p_{\varrho})^{m}$ ersetzen durch ein Produkt von $m$ Linearfaktoren $(A^{\varrho} | p_{\varrho})(B^{\varrho} | p_{\varrho}) \cdots (M^{\varrho} | p_{\varrho})$; wodurch die invarianten Bildungen übergehen in solche eines Systems von Linearformen, die erst nachträglich einander gleich gesetzt zu werden brauchen. Für jede einzelne Linearform aber sind die Bedingungen (12.), die bei Einführung von Differentialsymbolen nur Differentialquotienten 2.\ Ordnung enthalten, stets erfüllt.

\section*{\S\ 7. Entwicklung der allgemeinen Formen nach Normalformen.}
\addcontentsline{toc}{section}{\S\ 7. Entwicklung der allgemeinen Formen nach Normalformen}

Wir werden im folgenden (\S\ 8) eine Haupteigenschaft des Faltungsprozesses (41.) ableiten, seine Zusammensetzbarkeit aus Grundfaltungen, bedürfen dazu aber der Übertragung eines wichtigen, von Herrn Mertens\footnote{F.\ Mertens, \textit{Über invariante Gebilde quaternärer Formen} (siehe Einleitung), zitiert als ``Mertens''.} im quaternären Gebiete gegebenen Satzes auf das $n$-näre Gebiet. Dieser Satz sagt aus, daß sich ein beliebiges endliches System von Formen ersetzen läßt durch ein äquivalentes System von Normalformen. Unter ``Normalform'' verstehen wir dabei --- in Verallgemeinerung der binären und ternären Definition\footnote{Study, a.\ a.\ O.\ S.~54.} --- eine Form, deren sämtliche in den Koeffizienten lineare invariante Bildungen identisch verschwinden mit Ausnahme der Grundform selbst. Nach Satz VI (\S\ 6) genügt also eine Normalform $g$ dem System von Differentialgleichungen:
\begin{equation}\tag{46}
(G_{n-\varrho-\lambda} \frac{\partial}{\partial s} | p_{\varrho+\lambda}) = 0, \qquad \lambda = \tau, \tau+1, \ldots, n-\varrho,
\end{equation}
wo die Reihen $q_{n-\varrho-\tau}$,$p_{\varrho+\tau}$ als willkürliche Größen zu betrachten sind.

Diese Differentialgleichungen stimmen im quaternären Gebiet --- unter Berücksichtigung von (39.) für $\varrho = \tau = 2$ --- mit den von Herrn Mertens ohne Einführung der willkürlichen Reihen $p$,$q$ gegebenen 10 Differentialgleichungen (a.\ a.\ O., Formel (28)) vollständig überein; ihre Kenntnis, zusammen mit Satz VI, erlaubt uns, das dortige Beweisverfahren fast wörtlich auf $n$ Variable zu übertragen. Hinzuzufügen ist einzig der mittels der Zerlegungsidentität (30.) erbrachte Nachweis, daß die konstruierten Formen Normalformen sind; im quaternären Gebiet ergibt sich dieser Nachweis noch ohne weitere Umformung. Es genügt deshalb auch eine kurze Darlegung des Beweises, und zwar der Einfachheit halber in dem für das folgende nur in Betracht kommenden Fall einer einzigen Grundform $f$ und ihrer Formenreihe (vgl.\ \S\ 6 Schluß), da für invariante Bildungen beliebigen Grades in den Koeffizienten beliebig vieler Grundformen der Beweis derselbe bleibt.

Der Grundgedanke des Beweises ist, durch Einführung der transformierten Koeffizienten eine Form mit $n-1$ Reihen $p_{\varrho}$ zu ersetzen durch Formen mit $n$ den $x$ kogredienten Reihen $\xi$, den Reihen der Transformationsgrößen. Aus diesen Formen leiten sich durch invariante Differentiationsprozesse direkt die gesuchten Normalformen ab. Die Entwicklung der allgemeinen Formen nach Normalformen wird vermittelt durch eine Reihenentwicklung für Formen mit $\varrho$ ($\varrho < n$) kogredienten Reihen $\xi$; invariante Differentiationsprozesse führen zu den Formen in den ursprünglichen Reihen $p_{\varrho}$ zurück.

Es seien $\xi^{(1)}$,$\ldots$,$\xi^{(n)}$ die Reihen der Transformationsgrößen, so daß man hat:
\begin{equation}\tag{47}
x_{\alpha} = \sum_{\beta} \xi^{(\beta)}_{\alpha} y_{\beta}, \qquad p_{\varrho} = [\xi^{(i_{1})} \cdots \xi^{(i_{\varrho})}]_{\varrho}.
\end{equation}
Die transformierten Koeffizienten der Formen $f_{1}$,$f_{2}$,$\ldots$ seien mit $(f)$,$(\varphi)$,$\ldots$ bezeichnet, und es sei
\begin{equation}\tag{48}
(f) = \sum_{k_{1} \geq \cdots \geq k_{n}} c_{k_{1} \cdots k_{n}} (f)_{k_{1} \cdots k_{n}}, \qquad (\varphi) = \sum \cdots
\end{equation}
Es wird dann:
\begin{multline*}\tag{49}
C_{k_{1} \cdots k_{n}}(f) = (S^{\varrho_{1}} | p_{\varrho_{1}})^{m_{1}} \cdots (S^{\varrho_{k}} | p_{\varrho_{k}})^{m_{k}} \\
= (S^{\varrho_{1}} | \xi^{(1)} \cdots \xi^{(\varrho_{1})})^{n_{1}} (S^{\varrho_{1}} | \xi^{(2)} \cdots \xi^{(\varrho_{1}+1)})^{n_{2}} \cdots (S^{\varrho_{k}} | \xi^{(n)} \cdots)^{n_{k}},
\end{multline*}
wo $n_{1} + n_{2} + \cdots + n_{k} = m_{1}$ usf.\ Aus jedem Koeffizienten $(f)$, für den
\begin{equation}\tag{50}
k_{1} \geq k_{2} \geq k_{3} \geq \cdots \geq k_{n},
\end{equation}
ergibt sich durch den die Reihen $\xi$ genau erschöpfenden Differentiationsprozeß:
\begin{equation}\tag{51}
\varphi = \left( \frac{\partial^{k_{1}-k_{2}}}{\partial \xi^{(1)}_{1} \cdots} \right) \left( \frac{\partial^{k_{2}-k_{3}}}{\partial \xi^{(2)}_{1} \cdots} \right) \cdots (f)_{k_{1} \cdots k_{n}}
\end{equation}
die Form
\begin{equation}\tag{52}
g = (S^{\varrho_{1}} | p_{\varrho_{1}})^{m_{1}} (S^{\varrho_{2}} | p_{\varrho_{2}})^{m_{2}} \cdots (S^{\varrho_{k}} | p_{\varrho_{k}})^{m_{k}}.
\end{equation}

Die Formen $g$ sind die gesuchten Normalformen. Sie haben nämlich die folgenden, die Normalformen definierenden Eigenschaften:

\textit{1.}\ Sie sind in den Koeffizienten lineare invariante Bildungen der Grundform $f$. Dies folgt daraus, daß sie aus $f$ durch die invarianten Differentiationsprozesse (51.) und (52.) hervorgehen. Sie lassen sich also (nach \S\ 6 Schluß) linear ausdrücken durch Polaren der endlichen Formenreihe von $f$. Man erhält diese Polaren, indem man die in $g$ eingehenden Variabeln $p_{\varrho}$ auffaßt als Koeffizienten einer in Linearformen zerfallenden Form mit den Variabelnreihen $v$:
\begin{equation*}
H = (V^{\varrho_{1}} | p_{\varrho_{1}}) (V^{\varrho_{2}} | p_{\varrho_{2}}) \cdots (V^{\varrho_{k}} | p_{\varrho_{k}}).
\end{equation*}
Die Simultaninvarianten der Formenreihen $f$ und $H$ ergeben alle in Betracht kommenden Polaren, da diese letzteren in den einzelnen Reihen $p_{\varrho}$ von gleicher Dimension sein müssen wie $g$ selbst.

\textit{2.}\ Sie genügen den Differentialgleichungen (46.). Wendet man nämlich auf $g$ den Differentialprozeß (45.) an, so kommt:
\begin{equation*}
g' = (G_{n-\varrho-\lambda} S^{\varrho}_{1} \cdots S^{\varrho}_{\lambda+1} | p_{\varrho+\lambda}).
\end{equation*}
Aus der Zerlegungsidentität (30.) folgt, wenn man die aus den Reihen $s$ gebildeten Reihen $S^{\varrho} = s^{(1)} \cdots s^{(\varrho)}$, $p_{\varrho} = [s^{(1)} \cdots s^{(\varrho)}]_{\varrho}$ mit den dortigen $q_{\varrho}$,$p_{\varrho}$ identifiziert:
\begin{equation*}
g' = \sum_{\beta=\varrho-\tau+\lambda}^{\varrho} \varepsilon (q_{\varrho} S^{\varrho} | s^{(\beta)} p_{\varrho-1}).
\end{equation*}
Es wird dabei:
\begin{equation*}
D_{\varrho} = \sum_{\alpha=1}^{n} \frac{\partial^{2}}{\partial x^{(1)}_{\alpha} \cdots} = 0,
\end{equation*}
wo $i_{1}$,$\ldots$,$i_{\varrho}$ irgendwelche $\varrho$, aus den Zahlen $1$ bis $\varrho$ herausgegriffene, von einander verschiedene Zahlen bedeuten. Da nun $\lambda > \varrho-\tau$, so folgt, daß sich unter diesen Zahlen $i_{1}$,$\ldots$,$i_{\varrho}$ notwendig solche zwischen $1$ bis $\tau$ befinden; damit verschwindet aber $s^{(1)} \cdots s^{(\tau)}$ identisch, folglich jedes einzelne Matrizenprodukt und damit auch $g'$.

\section*{\S\ 8. Zurückführung des Faltungsprozesses auf Grundfaltungen.}
\addcontentsline{toc}{section}{\S\ 8. Zurückführung des Faltungsprozesses auf Grundfaltungen}

Im Anschluß an Satz VII führen wir jetzt die am Anfang vom \S\ 7 erwähnte Zurückführung des Faltungsprozesses auf Grundfaltungen durch. Wir bezeichnen dabei (vgl.\ \S\ 5) jede Faltung zweier kogredienter Reihen $S^{\varrho}$,$T^{\varrho}$ als ``kogrediente Faltung'', und speziell die kogrediente Faltung vom Defekt~$1$ zweier Reihen $S^{\varrho}$,$T^{\varrho}$ als ``Grundfaltung vom Typus $\varrho$''. Dann beweisen wir, in Vervollständigung von Satz VI, den

\medskip
\textbf{Satz VIII:} \textit{Der allgemeine Faltungsprozeß (41.) läßt sich zusammensetzen aus den $(n-1)$ Grundfaltungen vom Typus $\varrho$, wo $\varrho = 1, 2, \ldots n-1$. Mit anderen Worten: Zur Erzeugung aller invarianten Bildungen genügt die sukzessive Anwendung der $n-1$ Grundfaltungen.}
\medskip

Im binären Gebiet fällt der Faltungsprozeß (41.) direkt mit der einzigen möglichen Grundfaltung zusammen; im ternären Gebiet habe ich Satz VIII in \S\ 1 meiner Dissertation (vgl.\ Einleitung) bewiesen. Es sind dort, wegen (44.), die allgemeinen kogredienten Faltungen noch identisch mit den Grundfaltungen, und es ist bemerkenswert, daß Satz VIII für $n$ Variable die Reduktion auf Grundfaltungen leistet, und nicht nur die viel schwächere auf kogrediente Faltungen. Der Beweis ist im Prinzip dem ternären nachgebildet, indem die allgemeinsten Faltungen aus den Grundfaltungen konstruiert werden, unter Anwendung der symbolischen Identitäten. Wir werden erzeugen:
\begin{enumerate}
\item beliebige Faltungen vom Defekt~$1$ durch kogrediente vom Defekt~$1$, d.\ h.\ durch Grundfaltungen (Analogon zum ternären Fall),
\item beliebige Faltungen vom Defekt $\lambda$ durch kogrediente vom Defekt $\lambda$ und durch beliebige Faltungen vom Defekt~$1$,
\item kogrediente Faltungen vom Defekt $\lambda$ durch kogrediente vom Defekt $\lambda-1$ und durch beliebige Faltungen vom Defekt~$1$.
\end{enumerate}

Dabei setzen wir, was wesentlich ist, nach Satz VII die beiden zu faltenden Formen $S$ und $T$ als Normalformen voraus. Es gelten also, nach (46.), die Relationen:
\begin{equation*}
(G_{n-\varrho-\tau} S^{\varrho}_{1} \cdots | p_{\varrho+\tau}) = 0, \qquad (G_{n-\varrho-\tau} T^{\varrho}_{1} \cdots | p_{\varrho+\tau}) = 0.
\end{equation*}

Es genügt weiter, nach \S\ 2 (Schluß) die Formen $S$ und $T$ in allen Variabelnreihen linear anzunehmen; d.\ h.\ die konstruierten Formen gehören der Formenreihe an:
\begin{equation*}
f = (S^{\varrho_{1}} | p_{\varrho_{1}}) \cdots (S^{\varrho_{k}} | p_{\varrho_{k}}) (T^{\tau_{1}} | p_{\tau_{1}}) \cdots (T^{\tau_{l}} | p_{\tau_{l}}).
\end{equation*}

\textit{1) Erzeugung von $(G_{n-\varrho} S^{\varrho} T^{\tau} | p_{\varrho+1} p_{\tau-1})$, wo $\varrho > \tau$, aus Grundfaltungen vom Typus $\tau+\sigma$; $\sigma = 0, 1, \ldots \varrho-\tau$.}

Aus der Grundfaltung vom Typus $1$
\begin{equation*}
(w S^{\varrho} | p_{\varrho}) (V^{1} | p_{1}) = \varepsilon (w S^{\varrho} V^{1} | p_{\varrho+1})
\end{equation*}
erhalten wir durch die Substitution
\begin{equation*}
w = T^{\tau}, \qquad p_{\varrho} = p_{\tau},
\end{equation*}
eine Form der Formenreihe $f$, mit drei Reihen vom oberen Gewichtsindex $\tau+1$:
\begin{equation}\tag{63}
(w S^{\varrho} | p_{\tau}) (V^{1} | p_{1}) = \varepsilon (w S^{\varrho} V^{1} | p_{\tau+1}).
\end{equation}

Eine Grundfaltung vom Typus $\tau+1$, angewandt auf (63.), ergibt:
\begin{equation*}
(G_{n} | R^{\tau+1} Q^{\varrho-1}) = (R^{\tau+1} S^{\varrho} | p_{\tau+1} p_{\varrho-1});
\end{equation*}
aus (31.) ergibt sich daraus, durch die Substitution $q_{n-\varrho-1} w = q_{n-\varrho-1}$, der Wert:
\begin{equation*}
w S^{\varrho} = \varepsilon \cdot \ldots
\end{equation*}
Der Ausdruck unter dem Summenzeichen verschwindet nach (62.) identisch; wir sind also zu einer Form gelangt:
\begin{equation*}
(w S^{\varrho} | p_{\tau+1}) (V^{1} | p_{1}) = \varepsilon (w S^{\varrho} V^{1} | p_{\tau+2}),
\end{equation*}
die aus (63.) durch Verwandlung von $\tau$ in $\tau+1$ hervorgeht. Die $(\varrho-\tau)$-fache Wiederholung des Prozesses führt also zu der gesuchten Faltung:
\begin{equation*}
(w S^{\varrho} | p_{\varrho}) (V^{1} | p_{1}) = \varepsilon (w S^{\varrho} V^{1} | p_{\varrho+1}).
\end{equation*}

Wir deuten den durchgeführten Prozeß kurz an durch die Formel:
\begin{equation}\tag{64}
(G_{n-\varrho} S^{\varrho} T^{\tau} | p_{\varrho+1} p_{\tau-1}) = \varepsilon (G_{n-\varrho-1} S^{\varrho} | p_{\varrho+1}) \cdots (T^{\tau} | p_{\tau});
\end{equation}
wir sehen, daß nur die Normalformeigenschaft von $S$, nicht aber die von $T$, benutzt wird. Ein zu (64.) dualistischer Prozeß läßt sich unter alleiniger Benutzung der Normalformeigenschaft von $T$ in umgekehrter Reihenfolge durchführen, nämlich der folgende:
\begin{equation}\tag{65}
(G_{n-\tau} T^{\tau} S^{\varrho} | p_{\tau+1} p_{\varrho-1}) = \varepsilon' (G_{n-\tau-1} T^{\tau} | p_{\tau+1}) \cdots (S^{\varrho} | p_{\varrho}),
\end{equation}
nach den Relationen:
\begin{align*}
(G_{n-\tau} T^{\tau} | p_{\tau+1}) &= \varepsilon (T^{\tau} R^{1} | p_{\tau-1} p_{1}), \\
(u_{n-\tau} R^{1} | T^{\tau}_{0} T^{\tau}_{1} p_{\tau-2}) &= \varepsilon (T^{\tau}_{0} | p_{\tau-1}) (R^{1} | p_{1}), \\
(G_{n} | \text{Typ.}) &= \varepsilon (\text{usw.}) (1, 2, \ldots).
\end{align*}

Es mag noch auf den Zusammenhang mit der in \S\ 6 (Schluß) angegebenen Abnahme des Gesamtgewichtes hingewiesen werden, die für Faltungen vom Defekt~$1$ den Wert $2(1+(\varrho-\tau))$, für Grundfaltungen den Wert $2 \cdot 1$ hat. Wenn also die Zusammensetzbarkeit aus Grundfaltungen überhaupt möglich ist, so sind notwendig $\varrho-\tau+1$ solcher dazu erforderlich, wie oben durchgeführt.

\textit{2) Erzeugung allgemeiner Faltungen aus kogredienten und Grundfaltungen.} Die Gewichtsabnahme beträgt für eine allgemeine Faltung (41.):
\begin{equation*}
2(\lambda^{2} + \lambda(\varrho-\tau)) = 2[\lambda^{2} + (\varrho-\tau)(1+(\lambda-1))],
\end{equation*}
ergibt also die Möglichkeit einer Zerlegung in eine kogrediente Faltung vom Defekt $\lambda$ (Gewichtsabnahme $2\lambda^{2}$) und $\varrho-\tau$ Faltungen vom Defekt~$1$ und der Differenz der Gewichtsindizes $\lambda-1$ (Gewichtsabnahme $2\{1+(\lambda-1)\}$). Diese Zerlegung stimmt im Falle $\lambda=1$ mit der unter 1) angegebenen überein; wir zeigen, daß sie sich in der Tat realisieren läßt.

Aus der kogrediente Faltung vom Defekt $\lambda$
\begin{equation*}
(R^{\lambda} | S^{\varrho} T^{\tau} p_{\varrho})
\end{equation*}
erhalten wir durch die Substitution: $R^{\lambda} = T^{\tau}_{n-\tau} p_{\varrho}$, eine Form mit den beiden Reihen vom oberen Gewichtsindex $\tau+\lambda$ und $\tau+1$:
\begin{equation}\tag{66}
(R^{\lambda} S^{\varrho} | p_{\tau+\lambda} p_{\varrho}) = \varepsilon (R^{\lambda} | p_{\tau+\lambda}) (S^{\varrho} | p_{\varrho}).
\end{equation}
Die Reihe mit dem niedrigeren Gewichtsindex $\tau+1$ gehört einer Normalform an; (66.) läßt also eine aus $\lambda$ Grundfaltungen zusammengesetzte Faltung vom Defekt~$1$ zu, in der Reihenfolge von (65):
\begin{equation*}
(R^{1} | S^{\varrho}) \to (S^{\varrho} | p_{\varrho}) \to \cdots
\end{equation*}
Dies ergibt, unter Berücksichtigung von (31.) und (62.):
\begin{equation*}
\text{Vers } [S^{\varrho} | p_{\varrho}] = \varepsilon (R^{1} S^{\varrho} | p_{\varrho+1}),
\end{equation*}
also eine Form, die aus (66.) durch Verwandlung von $\tau$ in $\tau+1$ hervorgeht, analog wie dies unter 1) mit (63.) der Fall war. Die $(\varrho-\tau)$-fache Wiederholung des Prozesses führt zu der gesuchten Faltung:
\begin{equation*}
(R^{\lambda} S^{\varrho} | p_{\varrho+1}) = \varepsilon (G_{n} | \ldots).
\end{equation*}
Bei dem ganzen Prozeß war allein die Normalformeigenschaft von $S$ benutzt; die Betrachtung von $T$ als Normalform führt auf den dualistischen Prozeß, mit vollständiger Umkehrung der Reihenfolge, nach den Formeln:
\begin{align*}
(G_{n-\varrho-\lambda} S^{\varrho} | p_{\varrho+\lambda} p_{\tau-\lambda}) &= (T^{\tau} R^{\lambda} | p_{\tau-2}) \to \text{usw.}, \quad R^{\lambda} = \text{usw.}, \\
(R^{\lambda} | P^{\varrho} p_{\varrho+\lambda} p_{\tau-1}) &= \varepsilon (T^{\tau} | p_{\tau}) \cdots, \\
(G_{n-\varrho} T^{\tau} | T^{\tau} R^{\lambda} p_{\varrho}) &= \varepsilon (G_{n-\varrho} S^{\varrho} | p_{\varrho+\lambda} p_{\tau-\lambda}) + \cdots
\end{align*}

\textit{3) Erzeugung kogredienter Faltungen vom Defekt $\lambda$ aus kogredienten vom Defekt $\lambda-1$ und Grundfaltungen.} Die Gewichtsabnahme für kogrediente Faltungen vom Defekt $\lambda$ beträgt: $2\lambda^{2} = 2((\lambda-1)^{2} + 2\lambda - 1)$, läßt also die Möglichkeit einer Zerlegung in eine kogrediente Faltung vom Defekt $\lambda-1$ und $2\lambda-1$ Grundfaltungen zu. Wir finden dieselbe am einfachsten, wenn wir erst speziell setzen: $B^{\tau} = A^{\varrho}$. Die einzig mögliche kogrediente Faltung vom Defekt $\lambda$ wird: $(G_{n-2\varrho} S^{\varrho} | T^{\varrho}_{\varrho})$; sie enthält eine Reihe $u$ und $x$ weniger als die Faltung vom Defekt $\lambda-1$ ebenderselben Reihen $(w S^{\varrho} | T^{\varrho} x)$. Wir setzen umgekehrt die Faltung vom Defekt $\lambda$ aus Grundfaltungen, die das Verschwinden einer Reihe $u$ und $x$ bewirken (Gewichtsabnahme $2(n-1) = 2(2\lambda-1)$) und zugleich eine neue Symbolreihe $R^{\lambda}$ erzeugen, und aus einer kogrediente Faltung vom Defekt $\lambda-1$ zusammen.

Die einfachste Faltung, die das leistet, ist die folgende vom Defekt~$1$:
\begin{equation*}
(s T^{\tau} | \sigma p_{\tau}) = \varepsilon (S^{\varrho} | R^{\varrho}_{\lambda-1} p_{\varrho}) = \varepsilon (R^{\lambda} | p_{\varrho}),
\end{equation*}
wo gesetzt ist:
\begin{equation*}
R^{\lambda} = s T^{\tau}, \quad s = S^{\varrho}, \quad c = S_{\varrho}, \quad R^{\varrho}_{\lambda-1} = \text{usw.}
\end{equation*}
Sie ist nach (65.) und (64.) zusammengesetzt aus den $2\lambda-1$ Grundfaltungen:
\begin{equation}\tag{67}
D_{\lambda} = (G_{n-\varrho-1} S^{\varrho} | p_{\varrho+1}), \quad (G_{n-\varrho-2} S^{\varrho} | p_{\varrho+2}), \quad \ldots, \quad (G_{n-\varrho-\lambda} S^{\varrho} | p_{\varrho+\lambda}).
\end{equation}

Eine Faltung vom Defekt $\lambda-1$ ergibt, unter Berücksichtigung von (21.) und (62.):
\begin{equation*}
(w S^{\varrho} | \sigma R^{\varrho}_{\lambda-1}) = \varepsilon (R^{\lambda-1} | S^{\varrho} \sigma) = \varepsilon (s T^{\tau} | S^{\varrho} x) = \varepsilon (S^{\varrho} | T^{\varrho}_{\varrho}).
\end{equation*}

Der allgemeine Fall, für zwei Reihen $S^{\varrho}$,$T^{\tau}$ und beliebiges $n$, läßt sich nun direkt aus dem speziellen übertragen. An Stelle von (67.) treten die folgenden $2\lambda-1$ Grundfaltungen:
\begin{equation}\tag{68}
\Pi_{\lambda} = (G_{n-\varrho-\lambda} S^{\varrho} | p_{\varrho+\lambda}), \quad (G_{n-\tau-\lambda} T^{\tau} | p_{\tau+\lambda}), \quad \ldots
\end{equation}

Die erste Hälfte der Faltungen ergibt die Form:
\begin{equation*}
(G_{n-\varrho-\lambda} S^{\varrho} | p_{\varrho+\lambda} p_{\tau-\lambda})
\end{equation*}
aus der durch die Substitutionen
\begin{equation*}
S^{\varrho} = y \cdot \sigma, \qquad T^{\tau} = w \cdot \tau, \qquad R^{\lambda} = \text{usw.}
\end{equation*}
und Anwendung der zweiten Hälfte der Faltungen hervorgeht:
\begin{equation*}
(G_{n-\varrho-\lambda} S^{\varrho} | R^{\lambda} p_{\varrho-\lambda}) \to \cdots
\end{equation*}

Substituiert man:
\begin{equation*}
Y^{n-\varrho-\lambda} = y^{\sigma} \cdot \sigma Y, \qquad \ldots
\end{equation*}
so kommt durch eine kogrediente Faltung vom Defekt $\lambda-1$:
\begin{equation}\tag{69}
(R^{\lambda} S^{\varrho} | Y^{\lambda} p_{\varrho-\lambda+\tau}) \to \cdots
\end{equation}

Wir zeigen, daß dies die gesuchte Faltung vom Defekt $\lambda$ ist. Substituiert man
\begin{equation*}
R^{\lambda} = \sigma p_{\varrho-1}, \qquad p_{\varrho} = \text{usw.}, \qquad R^{\lambda-1}_{\varrho} = \text{usw.},
\end{equation*}
und beachtet, daß durch Auflösen von $y$
\begin{equation*}
M^{\varrho+1} = S^{\varrho} Y^{\lambda} - \varepsilon (G_{n} | \text{usw.}) = 0
\end{equation*}
kommt, so ergibt sich nach (20.) für (69.) der Wert:
\begin{equation*}
\varepsilon (S^{\varrho} R^{\lambda} | p_{\varrho+\lambda} y) = \varepsilon (Y^{n-\varrho-\lambda} D_{\varrho} | R^{\lambda}_{\varrho} \sigma^{\lambda-1}).
\end{equation*}

Diese Faltung ist aber vom Typus (43.), also äquivalent
\begin{equation*}
(A^{\varrho} | R^{\varrho}_{\lambda} p_{\varrho-\lambda}) = \varepsilon (W T^{\tau} | R^{\lambda}_{\varrho} p_{\varrho-\lambda}), \quad \text{wo } R^{\lambda}_{\varrho} = \text{usw.}
\end{equation*}
Hier ist $R^{\lambda}_{\varrho}$ schon von der gewünschten Schlußform; der Ausdruck entspricht der Form $(s T^{\tau} | S^{\varrho} x)$ im Spezialfall. Beachtet man nun, daß durch Auflösen von $w$ und $R^{\lambda}_{\varrho}$ sich ergibt:
\begin{align*}
(w R^{\lambda} | R^{\lambda}_{\varrho} S^{\varrho}) &= \varepsilon (H^{\lambda}_{\varrho+\tau} | S^{\varrho} \ldots) \\
&= \varepsilon (G_{n-\varrho} S^{\varrho} | D_{\varrho} \ldots) = 0,
\end{align*}
so kommt nach (21.) der Wert:
\begin{equation*}
(w q_{\varrho-1} | \ldots R^{\lambda}) = \varepsilon (A^{\varrho} | R^{\lambda}_{\varrho} [S^{\varrho} \sigma y \ldots y^{\varrho}])
\end{equation*}
äquivalent $(q_{\varrho-\lambda-1} S^{\varrho} | T^{\varrho}_{\varrho-\lambda} p_{\varrho-\lambda})$.

Bei dem ganzen Prozeß war nur die Normalformeigenschaft von $S$, nicht die von $T$ benutzt; die kogrediente Faltung vom Defekt $\lambda-1$ zur Erzeugung von (69.) läßt sich also ersetzen durch $2\lambda-3$ Grundfaltungen und eine kogrediente Faltung vom Defekt $\lambda-2$, die ihrerseits wieder eine Zerlegung in Grundfaltungen zuläßt usf. Analog ergibt sich die Erzeugung bei Benutzung der Normalformeigenschaft von $T$ und Anwendung des dualistischen Prozesses, d.\ h.\ der Grundfaltungen:
\begin{equation*}
\Pi^{*}_{\lambda} = (G_{n-\tau-\lambda} T^{\tau} | p_{\tau+\lambda}), \quad (G_{n-\varrho-\lambda} S^{\varrho} | p_{\varrho+\lambda}), \quad \ldots
\end{equation*}

Damit ist Satz VIII vollständig bewiesen.

\end{document}
